%=====================================================================
%  Design of a Three-Phase N-Channel MOSFET Inverter and Analog
%  Front-End for a 2.5-5 kW Battery-Powered BLDC/PMSM Drive with
%  Field-Oriented Control.
%
%  Combined English design report (IEEEtran). Compile with pdflatex
%  (pdflatex -> pdflatex). Requires IEEEtran.cls, circuitikz, siunitx.
%=====================================================================
\documentclass[journal,onecolumn,11pt]{IEEEtran}

\usepackage[T1]{fontenc}
\usepackage{amsmath,amssymb}
\usepackage{mathtools}
\usepackage{siunitx}
\usepackage{booktabs}
\usepackage{array}
\usepackage{multirow}
\usepackage{tabularx}
\usepackage{graphicx}
\usepackage[siunitx]{circuitikz}
\usepackage{tikz}
\usetikzlibrary{arrows.meta,positioning,calc,shapes.geometric}
\usepackage{enumitem}
\setlist{itemsep=2pt,topsep=3pt}
\usepackage[hidelinks,colorlinks=true,linkcolor=blue,citecolor=blue,urlcolor=blue]{hyperref}
\usepackage{url}
\urlstyle{sf}

\usepackage{textcomp}
\sisetup{detect-all,per-mode=symbol,group-digits=integer}
% Default \ohm glyph is missing in the pdflatex Times TS1 encoding and prints as
% a tofu box; force the unit to use the (always-available) math capital omega.
\DeclareSIUnit{\ohm}{\ensuremath{\Omega}}

% Convenience macros
\newcommand{\Vbus}{V_{\mathrm{bus}}}
\newcommand{\Rshunt}{R_{\mathrm{sh}}}
\newcommand{\Iph}{I_{\mathrm{ph}}}
\newcommand{\fsw}{f_{\mathrm{sw}}}
\newcommand{\Rdson}{R_{DS(\mathrm{on})}}

% Left-aligned wrapping column for tabularx
\newcolumntype{L}{>{\raggedright\arraybackslash}X}

\begin{document}

\title{Design of a Three-Phase N-Channel MOSFET Inverter and Analog Front-End
for a \mbox{2.5--5\,kW} Battery-Powered BLDC/PMSM Drive with Field-Oriented Control}

\author{Hardware Design Report\\[2pt]
\small Power stage: Infineon IPT015N10N5 (\SI{100}{\volt}, \SI{1.5}{\milli\ohm})
OptiMOS\textsuperscript{\texttrademark}\,5 MOSFET $+$ Texas Instruments UCC27211A-Q1 half-bridge gate driver\\
\small Controller: STMicroelectronics STM32G473RCT6 \quad
Switching frequency: \SI{40}{\kilo\hertz} \quad Control: field-oriented control (FOC) with space-vector PWM}

\markboth{Hardware Design Report: Three-Phase BLDC/PMSM Inverter and Analog Front-End}{}
\maketitle

%=====================================================================
\begin{abstract}
This report presents the component-level design of the power stage, the gate-drive stage, and
the analog front-end of a three-phase voltage-source inverter for a brushless DC / permanent-magnet
synchronous (BLDC/PMSM) motor under field-oriented control (FOC) with space-vector pulse-width
modulation (SVPWM). The drive is battery-powered: a \SI{24}{\volt} pack whose terminal voltage
ranges from \SI{19}{\volt} to \SI{30}{\volt} with state of charge, delivering up to \SI{2.5}{\kilo\watt};
and a \SI{48}{\volt} pack ranging from \SI{39}{\volt} to \SI{60}{\volt}, delivering up to \SI{5}{\kilo\watt}.
The variation of the pack voltage is central to the design: the \emph{highest} bus voltage
(\SI{60}{\volt}) sets the blocking and $\mathrm{d}v/\mathrm{d}t$ stresses, while the \emph{lowest}
bus voltage (\SI{19}{\volt}) produces the largest current for a given power and therefore the
worst case for conduction loss and for the current-sense chain. Each switch position uses two
Infineon IPT015N10N5 \SI{100}{\volt} OptiMOS\textsuperscript{\texttrademark}\,5 MOSFETs in
parallel (four devices per half-bridge leg, twelve total), and each leg is driven by one Texas
Instruments UCC27211A-Q1 \SI{120}{\volt}, \SI{3.7}{\ampere}/\SI{4.5}{\ampere} half-bridge gate
driver supplied at \SI{12}{\volt}. The report (i) verifies the electrical compatibility of the
selected devices across the full battery-voltage range; (ii) derives the gate-drive passive
network --- bootstrap capacitor, bootstrap-diode requirement, asymmetric gate resistors,
gate--source pull-downs, and decoupling --- from published design equations; and (iii) designs
the three remaining analog blocks of the input stage: the in-line shunt phase-current-sense
chain, the DC-bus and phase voltage-sense chains, and a do-not-populate RC snubber. The
worst-case continuous phase current is derived as $\approx\SI{119}{\ampere}$ RMS
($\approx\SI{169}{\ampere}$ peak) and used to size a \SI{0.1}{\milli\ohm} metal-element shunt read
by a Texas Instruments INA240A2 current-sense amplifier with enhanced PWM rejection. Every
quantitative claim and every design equation is supported by a primary source.
\end{abstract}

\begin{IEEEkeywords}
BLDC, PMSM, field-oriented control (FOC), space-vector PWM, MOSFET gate driver, bootstrap,
asymmetric gate resistor, in-line current sensing, shunt resistor, INA240, voltage divider,
STM32G4, RC snubber, single ground.
\end{IEEEkeywords}

%=====================================================================
\section{Introduction and Scope}
\IEEEPARstart{T}{he} target system is a non-isolated three-phase voltage-source inverter, built
on a single shared ground, intended to drive any BLDC/PMSM motor of roughly \SI{0.5}{} to
\SI{5}{\kilo\watt} under field-oriented control with space-vector PWM and Hall-effect rotor-position
sensing. The controller is an STMicroelectronics STM32G473RCT6 --- a motor-control microcontroller
based on an Arm\textsuperscript{\textregistered} Cortex\textsuperscript{\textregistered}-M4 core
that integrates up to five \mbox{12-bit} analog-to-digital converters (ADCs), operational
amplifiers usable as a programmable-gain amplifier (PGA) or unity-gain follower with their outputs
routed internally to an ADC, analog comparators, advanced motor-control timers, and a voltage
reference buffer (VREFBUF)~\cite{stm32g473ds}. The gate driver is the Texas Instruments
UCC27211A-Q1~\cite{ucc27211}, and each switch is an Infineon OptiMOS\textsuperscript{\texttrademark}\,5
IPT015N10N5 (\SI{100}{\volt}, \SI{1.5}{\milli\ohm})~\cite{ipt015ds,ipt015part}. The microcontroller,
the gate driver, and the choice of power MOSFET are fixed inputs to this work.

The report addresses, in one place, both the power/gate-drive stage and the analog input stage:
\begin{enumerate}[leftmargin=1.5em]
  \item the power switch and gate driver, their key parameters, and a compatibility assessment
        across the full battery-voltage range (Sections~\ref{sec:spec}--\ref{sec:compat});
  \item the gate-drive passive network, derived quantitatively (Section~\ref{sec:passive});
  \item the \emph{phase-current measurement} chain (in-line shunt $+$ amplifier $+$ ADC interface)
        for all three phases (Section~\ref{sec:current});
  \item the \emph{voltage measurement} chains for the DC bus and for each of the three phase
        nodes (Section~\ref{sec:voltage}); and
  \item an \emph{RC snubber} across the power switches (Section~\ref{sec:snubber}),
\end{enumerate}
followed by the grounding/noise strategy, thermal considerations, and the bill of materials. For
each analog block the report explains the operating principle, derives the governing equations,
computes numerical values for the present specification, and selects a specific, orderable
component.

Because the inverter accepts any \SI{24}{\volt} or \SI{48}{\volt} battery with no assumed
chemistry, the design is dimensioned for the worst case of each electrical stress rather than for
a single nominal operating point. The system uses a \emph{single shared ground}; noise immunity is
obtained through layout discipline rather than through a separate analog ground plane
(Section~\ref{sec:noise}).

%=====================================================================
\section{Design Specification}\label{sec:spec}

Table~\ref{tab:spec} lists the given design parameters. The control method is field-oriented
control with space-vector PWM throughout; six-step (trapezoidal) commutation is discussed only as
historical background in Section~\ref{sec:why}. Rotor position is obtained from Hall-effect
sensors.

\begin{table}[!t]
\centering
\caption{Given design parameters.}
\label{tab:spec}
\renewcommand{\arraystretch}{1.25}
\begin{tabularx}{\textwidth}{@{}>{\raggedright\arraybackslash}p{3.4cm} >{\raggedright\arraybackslash}p{4.0cm} L@{}}
\toprule
\textbf{Parameter} & \textbf{Value} & \textbf{Rationale / source} \\
\midrule
Power source & Battery (\SI{24}{\volt} or \SI{48}{\volt} pack) & Bus voltage varies with state of charge \\
DC-bus voltage range & \SI{24}{\volt} pack: \SIrange{19}{30}{\volt}\newline \SI{48}{\volt} pack: \SIrange{39}{60}{\volt} & $\Vbus^{\max}=\SI{60}{\volt}$ (full \SI{48}{\volt} pack); $\Vbus^{\min}=\SI{19}{\volt}$ (depleted \SI{24}{\volt} pack) \\
Output power & $\approx\SI{2.5}{\kilo\watt}$ (\SI{24}{\volt} pack) / $\approx\SI{5}{\kilo\watt}$ (\SI{48}{\volt} pack) & Motor power scales with bus voltage \\
Switching frequency $\fsw$ & \SI{40}{\kilo\hertz} & Above the audible band ($\sim\SI{20}{\kilo\hertz}$), so the PWM tone is inaudible \\
Control & FOC with SVPWM & Sinusoidal, rotor-field-oriented current control~\cite{ti_foc,ti_svpwm} \\
Position sensing & Hall-effect sensors & Six \ang{60} sectors per electrical revolution; FOC angle by interpolation \\
Power switch & IPT015N10N5, $\times 2$ in parallel per position & \SI{100}{\volt}, \SI{1.5}{\milli\ohm} OptiMOS\textsuperscript{\texttrademark}\,5, optimized for high-current applications~\cite{ipt015ds,ipt015part} \\
Gate driver & UCC27211A-Q1, $\times 1$ per leg & \SI{120}{\volt} boot, \SI{3.7}{\ampere}/\SI{4.5}{\ampere} half-bridge driver with integrated boot diode~\cite{ucc27211,ucc27211folder} \\
Gate-drive supply $V_{DD}$ & \SI{12}{\volt} & Within the device operating range \SIrange{8}{17}{\volt}~\cite{ucc27211} \\
\bottomrule
\end{tabularx}
\end{table}

\subsection{Worst-case phase current}\label{sec:iph}
The current-handling hardware must be dimensioned for the largest phase current the inverter can
carry, since this sets both the conduction loss in the MOSFETs and the shunt/amplifier scaling.
For a sinusoidally modulated three-phase inverter using SVPWM, the maximum fundamental phase-voltage
amplitude is $\hat V_{\mathrm{ph}}=\Vbus/\sqrt{3}$, so the maximum RMS phase voltage is
\begin{equation}
V_{\mathrm{ph,rms}}=\frac{\Vbus}{\sqrt{6}} .
\label{eq:vph}
\end{equation}
With three phases and displacement power factor $\mathrm{PF}=\cos\varphi$, the AC power is
$P = 3\,V_{\mathrm{ph,rms}}\,\Iph\,\mathrm{PF}$. Taking converter efficiency $\eta\!\to\!1$ (the
conservative assumption that maximizes current) and solving for the RMS phase current,
\begin{equation}
\boxed{\;\Iph = \frac{\sqrt{6}}{3}\,\frac{P}{\Vbus\,\mathrm{PF}}
        \approx 0.816\,\frac{P}{\Vbus\,\mathrm{PF}}\;}
\label{eq:iph}
\end{equation}
which is largest at minimum bus voltage and minimum power factor. Table~\ref{tab:env} evaluates
\eqref{eq:iph} at the four supply/power corners ($\mathrm{PF}=0.9$).

\begin{table}[!t]
\centering
\caption{Operating corners and resulting phase current from \eqref{eq:iph} ($\mathrm{PF}=0.9$).}
\label{tab:env}
\renewcommand{\arraystretch}{1.2}
\begin{tabular}{l c c c c}
\toprule
Supply / motor & $\Vbus$ (V) & $P$ (kW) & $\Iph$ (A,rms) & $\hat I_{\mathrm{ph}}$ (A,pk)\\
\midrule
\SI{24}{\volt} pack, \SI{2.5}{\kilo\watt} & 19 (min) & 2.5 & \textbf{119.4} & \textbf{168.8}\\
\SI{24}{\volt} pack, \SI{2.5}{\kilo\watt} & 30 (max) & 2.5 & 75.6 & 106.9\\
\SI{48}{\volt} pack, \SI{5}{\kilo\watt}   & 39 (min) & 5.0 & 116.3 & 164.5\\
\SI{48}{\volt} pack, \SI{5}{\kilo\watt}   & 60 (max) & 5.0 & 75.6 & 106.9\\
\bottomrule
\end{tabular}
\end{table}

The limiting corner is \SI{2.5}{\kilo\watt} at \SI{19}{\volt}: $\Iph^{\max}\approx\SI{119}{\ampere}$
RMS, $\hat I_{\mathrm{ph}}^{\max}\approx\SI{169}{\ampere}$ peak. The corresponding worst-case DC-bus
current is $P/\Vbus^{\min}=2500/19\approx\SI{132}{\ampere}$. Because two devices are paralleled in
each switch position, each MOSFET carries approximately half of the phase current, i.e.\
$\approx\SI{60}{\ampere}$ RMS, which is far below the $2\times\SI{250}{\ampere}$ ($T_C=\SI{100}{\celsius}$)
continuous rating of the paralleled pair~\cite{ipt015ds}. The per-device conduction loss at this
corner is $P=I^2\Rdson\approx 60^2\times\SI{1.5}{\milli\ohm}\approx\SI{5.4}{\watt}$ (rising toward
$\sim\SI{8}{\watt}$ as $\Rdson$ increases with junction temperature), which requires adequate
cooling (Section~\ref{sec:thermal}). All current-sense quantities below refer to this limiting
corner.

%=====================================================================
\section{Background and Theory of Operation}\label{sec:why}

\subsection{Why an inverter is required to drive a BLDC/PMSM machine}
Despite its name, the brushless DC motor is a synchronous AC machine: it has a permanent-magnet
rotor and a three-phase stator winding, and it lacks the mechanical commutator and brushes that a
conventional DC motor uses to switch current between windings as the rotor turns~\cite{ti_foc}.
The commutation function once performed mechanically by the brushes must therefore be performed
electronically. The inverter is the circuit that does this: it converts the fixed-polarity DC bus
voltage (here supplied by a battery) into the set of controlled, alternating phase currents needed
to keep the stator magnetic field rotating ahead of the rotor and thereby produce continuous
torque~\cite{ti_foc}.

This design uses \emph{field-oriented control} (FOC) with \emph{space-vector PWM} (SVPWM). In FOC
the measured three-phase currents are transformed (Clarke and Park transformations) into a
rotor-aligned reference frame, where a direct-axis ($d$) and a quadrature-axis ($q$) current are
regulated independently; the $q$-axis current produces torque while the $d$-axis current is
normally held at zero~\cite{ti_foc}. The resulting voltage demand is synthesized by SVPWM, which
selects and time-weights the inverter's switching states to approximate the commanded voltage
vector with the best use of the available bus voltage~\cite{ti_svpwm}. Compared with the simpler
six-step (trapezoidal, \ang{120}) scheme --- in which the six switches step through six discrete
field orientations per electrical revolution --- FOC drives sinusoidal phase currents and produces
smooth torque with low ripple, which is why it is adopted here. The rotor angle that the Park
transformation and SVPWM require is obtained from the Hall-effect sensors, interpolated between the
six \ang{60} sector transitions to recover a continuous electrical angle.

\subsection{How the inverter circuit works}
The power stage is a three-phase bridge: three identical half-bridge ``legs,'' one per motor phase,
each consisting of a high-side switch and a low-side switch in series across the DC bus, with the
motor phase connected to the midpoint. In this design each switch is two IPT015N10N5 MOSFETs in
parallel, so one leg is four MOSFETs and the full bridge is twelve.

Within a leg, the two switches operate complementarily --- never both on at once, which would
short-circuit the DC bus (``shoot-through''). Turning on the high-side switch connects the phase to
the positive bus; turning on the low-side switch connects it to ground; turning both off lets the
phase float. Modulating the three legs with SVPWM circulates sinusoidal currents through the phases
in the sequence that produces a rotating stator field~\cite{ti_svpwm}. The MOSFET body diodes
(together with the motor inductance) provide the freewheeling path for the inductive phase current
during the dead-time between switch states; the IPT015N10N5 body diode is characterized for this
purpose ($V_{SD}\approx\SI{0.85}{\volt}$, $Q_{rr}\approx\SI{316}{\nano\coulomb}$)~\cite{ipt015ds}.
The resistive on-state of the MOSFET ($\Rdson$ as low as \SI{1.5}{\milli\ohm}, again halved by
paralleling) keeps the conduction loss low at currents on the order of \SI{100}{\ampere}~\cite{ipt015ds,balogh}.

\subsection{Purpose of each component}
\begin{itemize}[leftmargin=1.4em]
  \item \textbf{Power MOSFETs (IPT015N10N5 $\times2$ per position):} the controlled switches that
        connect each phase to the buses. Two in parallel halve the effective on-resistance and
        share the conduction loss and heat; the positive temperature coefficient of $\Rdson$ makes
        the sharing self-stabilizing~\cite{balogh}.
  \item \textbf{Gate driver (UCC27211A-Q1):} converts the controller's PWM logic into the
        high-current \SI{12}{\volt} gate drive needed to switch the MOSFETs quickly, level-shifts
        the high-side signal to the floating switch node, and enforces under-voltage lockout
        (UVLO)~\cite{ucc27211}. A gate driver is required because a logic output (often
        \SI{3.3}{\volt}, $\le\SI{1}{\ampere}$) can neither fully enhance a power MOSFET nor supply
        the peak current for fast switching, and because the high-side gate must be referenced to a
        node that swings to the bus voltage~\cite{ucc27211,balogh}.
  \item \textbf{Bootstrap capacitor:} the floating energy reservoir that supplies the high-side
        driver while the switch node is high (Section~\ref{sec:boottheory}).
  \item \textbf{Bootstrap diode (internal):} a one-way valve that refills the bootstrap capacitor
        from $V_{DD}$ when the switch node is low and blocks the bus voltage when it is
        high~\cite{ucc27211,ir_an978}.
  \item \textbf{Gate resistors:} set and limit the gate current, shape $\mathrm{d}v/\mathrm{d}t$
        for EMI, and damp/balance the paralleled gates~\cite{balogh}.
  \item \textbf{Gate--source pull-down resistors:} hold the MOSFETs off at power-up and add
        $\mathrm{d}v/\mathrm{d}t$ immunity against false turn-on~\cite{balogh}.
  \item \textbf{Decoupling capacitors:} supply the fast gate-current pulses locally so that the
        \SI{12}{\volt} rail and the driver bias do not collapse during switching~\cite{ucc27211}.
  \item \textbf{DC-link capacitance:} stiffens the bus against the large half-bridge ripple current
        and limits the commutation overshoot within the \SI{100}{\volt} device rating.
\end{itemize}

%=====================================================================
\section{Power Switch: Infineon IPT015N10N5}\label{sec:fet}

The IPT015N10N5 is an N-channel OptiMOS\textsuperscript{\texttrademark}\,5 power transistor in a
TO-Leadless (PG-HSOF-8, ``TOLL'') package, which Infineon describes as suited to high-frequency
switching and intended for light-electric-vehicle, forklift, and power-tool drives~\cite{ipt015ds,ipt015part}.
The parameters relevant to this design are summarized in Table~\ref{tab:fet}; all values are from
the datasheet (Rev.~2.4) unless noted~\cite{ipt015ds}.

\begin{table}[!t]
\centering
\caption{Key parameters of the IPT015N10N5~\cite{ipt015ds}.}
\label{tab:fet}
\renewcommand{\arraystretch}{1.2}
\begin{tabularx}{\textwidth}{@{}l L L@{}}
\toprule
\textbf{Symbol} & \textbf{Parameter} & \textbf{Value / note} \\
\midrule
$V_{DS}$ & Drain--source voltage & \SI{100}{\volt} (absolute maximum) \\
$\Rdson$ & On-resistance at $V_{GS}=\SI{10}{\volt}$ & \SI{1.3}{\milli\ohm} typ.\ / \SI{1.5}{\milli\ohm} max. \\
$I_D$ & Continuous drain current & \SI{353}{\ampere} ($T_C=\SI{25}{\celsius}$); \SI{250}{\ampere} ($T_C=\SI{100}{\celsius}$) \\
$V_{GS(\mathrm{th})}$ & Gate threshold voltage & \num{2.2} / \num{3.0} / \SI{3.8}{\volt} (min/typ/max) \\
$V_{\mathrm{plateau}}$ & Miller-plateau voltage & \SI{4.4}{\volt} typ.\ (at $I_D=\SI{100}{\ampere}$) \\
$C_{iss}$ & Input capacitance & \SI{12}{\nano\farad} typ.\ / \SI{16}{\nano\farad} max. \\
$C_{oss}$ & Output capacitance & \SI{1.8}{\nano\farad} typ. \\
$C_{rss}$ & Reverse-transfer capacitance & \SI{80}{\pico\farad} typ.\ / \SI{140}{\pico\farad} max. \\
$R_{G,i}$ & Internal gate resistance & \SI{1.4}{\ohm} typ.\ / \SI{2.1}{\ohm} max. \\
$Q_{g}$ & Total gate charge (0--\SI{10}{\volt}) & \SI{169}{\nano\coulomb} typ.\ / \SI{211}{\nano\coulomb} max. \\
$Q_{gs}$ & Gate--source charge & \SI{53}{\nano\coulomb} \\
$Q_{gd}$ & Gate--drain (Miller) charge & \SI{34}{\nano\coulomb} typ.\ / \SI{51}{\nano\coulomb} max. \\
$g_{fs}$ & Transconductance & \SI{140}{\siemens} min.\ / \SI{280}{\siemens} typ. \\
$V_{SD}$ & Body-diode forward voltage & \SI{0.85}{\volt} typ.\ / \SI{1.2}{\volt} max. \\
$Q_{rr}$ & Body-diode reverse-recovery charge & $\approx\SI{316}{\nano\coulomb}$ typ. \\
$V_{GS,\max}$ & Gate--source voltage limit & $\pm\SI{20}{\volt}$ \\
\bottomrule
\end{tabularx}
\end{table}

A property that directly justifies paralleling is the positive temperature coefficient of $\Rdson$
(about $+0.7$ to $+\SI{1}{\percent\per\kelvin}$ in power MOSFETs): a device that carries more current
heats up, its resistance rises, and current is pushed back toward the cooler device, so paralleled
MOSFETs self-balance their currents~\cite{balogh}. This is the standard reason MOSFETs (unlike
bipolar transistors) are well suited to parallel operation~\cite{balogh}.

%=====================================================================
\section{Gate Driver: Texas Instruments UCC27211A-Q1}\label{sec:gd}

The UCC27211A-Q1 is an AEC-Q100-qualified \SI{120}{\volt} high-side/low-side half-bridge gate
driver that drives two N-channel MOSFETs from independent inputs; its floating high-side stage can
operate to \SI{120}{\volt}, and a \SI{120}{\volt} bootstrap diode is integrated on-chip, so no
external boot diode is mandatory~\cite{ucc27211,ucc27211folder}. The relevant parameters, all from
datasheet SLUSCG0C, are listed in Table~\ref{tab:gd}~\cite{ucc27211}.

\begin{table}[!t]
\centering
\caption{Key parameters of the UCC27211A-Q1~\cite{ucc27211}.}
\label{tab:gd}
\renewcommand{\arraystretch}{1.2}
\begin{tabularx}{\textwidth}{@{}l L L@{}}
\toprule
\textbf{Symbol} & \textbf{Parameter} & \textbf{Value} \\
\midrule
$V_{DD}$ & Supply (operating / abs.\ max) & \SIrange{8}{17}{\volt} / \SI{20}{\volt} \\
$I_{O,\mathrm{source}}$ / $I_{O,\mathrm{sink}}$ & Peak output current & \SI{3.7}{\ampere} / \SI{4.5}{\ampere} \\
$V_{DD,\mathrm{UVLO}}$ & UVLO threshold on $V_{DD}$ (rising) & \SI{7.0}{\volt} (hysteresis \SI{0.5}{\volt}) \\
$V_{HB,\mathrm{UVLO}}$ & UVLO threshold on high side (HB--HS, rising) & \SI{6.7}{\volt} (hysteresis \SI{1.1}{\volt}) \\
$V_{HB}$ & Max boot-supply voltage & \SI{120}{\volt} (HB--HS $\le\SI{20}{\volt}$) \\
$V_{HS}$ & Switch-node range & $-(24-V_{DD})\,$V (i.e.\ $-\SI{12}{\volt}$ at $V_{DD}=\SI{12}{\volt}$) to \SI{105}{\volt} recommended \\
$V_{F}$ & Boot-diode forward drop & \SI{0.85}{\volt} (@\SI{100}{\micro\ampere}) / \SI{1.05}{\volt} (@\SI{100}{\milli\ampere}) max. \\
$R_{D}$ & Boot-diode dynamic resistance & \SI{0.55}{\ohm} typ. \\
$t_{PD}$ & Propagation delay & \SI{20}{\nano\second} typ. \\
$t_{R}/t_{F}$ & Rise / fall time (\SI{1}{\nano\farad}) & \SI{7.2}{\nano\second} / \SI{5.5}{\nano\second} \\
$I_{HB}$ & Boot quiescent current & \SI{0.12}{\milli\ampere} max. \\
$SR_{HS}$ & Max high-side slew rate & \SI{50}{\volt\per\nano\second} \\
$C_{HB}$ & Recommended boot capacitor & \SIrange{0.022}{0.1}{\micro\farad} (gate-charge dependent) \\
$C_{VDD}$ & Recommended $V_{DD}$ decoupling & \SIrange{0.22}{4.7}{\micro\farad} \\
\bottomrule
\end{tabularx}
\end{table}

TI divides the device into five functional blocks --- input stages, UVLO protection, level shift,
boot diode, and output stages~\cite{ucc27211}. The inputs are independent HI/LI with
$\approx\SI{68}{\kilo\ohm}$ pull-downs and TTL-compatible thresholds with wide hysteresis for noise
immunity; they tolerate $-\SI{10}{\volt}$ to $+\SI{20}{\volt}$ and are independent of $V_{DD}$, so
they accept both microcontroller logic and analog-controller signals~\cite{ucc27211}. Separate
monitors on $V_{DD}$ and on the floating $V_{HB}$--$V_{HS}$ supply hold the outputs low until the
bias is valid. The level shifter translates the ground-referenced high-side command up to the
floating high-side stage with tight ($\approx\SI{4}{\nano\second}$) channel-to-channel delay
matching. Both output stages provide \SI{3.7}{\ampere} source / \SI{4.5}{\ampere} sink with fast,
low-impedance edges and a \SI{20}{\nano\second} propagation delay~\cite{ucc27211}.

The largest output current is required \emph{around the Miller-plateau voltage} ($\sim\SI{5}{\volt}$),
because it is there that the gate driver must push the Miller charge $Q_{gd}$ through the plateau in
the intended transition time; the datasheet states the peak-current requirement as
$I_{\mathrm{PEAK}}=Q_{GD}/\mathrm{d}t$, and the \SI{3.7}{\ampere} capability is checked against this
in Section~\ref{sec:gateR}~\cite{ucc27211,balogh}.

%=====================================================================
\section{Compatibility Assessment}\label{sec:compat}

Table~\ref{tab:compat} checks the device pair against the full battery-voltage range.

\begin{table}[!t]
\centering
\caption{Component compatibility check.}
\label{tab:compat}
\renewcommand{\arraystretch}{1.25}
\begin{tabularx}{\textwidth}{@{}L L L >{\raggedright\arraybackslash}p{2.3cm}@{}}
\toprule
\textbf{Check} & \textbf{Requirement} & \textbf{This design} & \textbf{Outcome} \\
\midrule
Blocking voltage & $V_{DS}\gg\Vbus^{\max}$ incl.\ ringing & \SI{100}{\volt} vs \SI{60}{\volt}; $\approx1.67\times$ DC margin & \textbf{Pass} --- $\approx\SI{40}{\volt}$ headroom for turn-off overshoot~\cite{ipt015ds} \\
High-side supply voltage & $V_{HB}>\Vbus^{\max}$ & \SI{120}{\volt} boot, \SI{105}{\volt} HS rec., vs \SI{60}{\volt} & \textbf{Pass}~\cite{ucc27211} \\
Gate-drive level & $V_{GS,\max}>V_{DD}\ge$ full enhancement & $\pm\SI{20}{\volt}$ limit, driven at \SI{12}{\volt} & \textbf{Pass}~\cite{ipt015ds,ucc27211} \\
UVLO & $V_{DD}>$ UVLO & \SI{12}{\volt} vs \SI{7}{\volt} ($V_{DD}$), \SI{6.7}{\volt} (HB) & \textbf{Pass}~\cite{ucc27211} \\
Negative node ringing & HS tolerates undershoot & $-\SI{12}{\volt}$ transient capability at $V_{DD}=\SI{12}{\volt}$ & \textbf{Pass}~\cite{ucc27211} \\
Current capability & $I_D\gg$ leg current & $2\times\SI{250}{\ampere}$ (\SI{100}{\celsius}) vs $\approx\SI{60}{\ampere}$/device worst case & \textbf{Pass}~\cite{ipt015ds} \\
Drive strength vs charge & Sufficient current through the Miller plateau for the \emph{paralleled} charge & see Section~\ref{sec:gateR} & \textbf{Pass} \\
\bottomrule
\end{tabularx}
\end{table}

The pair is electrically sound across the full battery range. At the highest bus voltage
(\SI{60}{\volt}, full \SI{48}{\volt} pack) the \SI{100}{\volt} class gives a $\approx1.67\times$ DC
margin and about \SI{40}{\volt} of reserve for the $L\,\mathrm{d}i/\mathrm{d}t$ turn-off overshoot
of $V_{DS}$; this makes control of the turn-off speed (via the asymmetric gate network of
Section~\ref{sec:asym}) more important than it would be on a fixed \SI{48}{\volt} bus. The only
point that genuinely required quantitative verification --- because the design parallels two
high-$Q_g$ MOSFETs per switch --- is whether one channel of the UCC27211A-Q1 can charge and
discharge the \emph{doubled} gate charge fast enough; this is treated in Section~\ref{sec:gateR}
and in the bootstrap sizing of Section~\ref{sec:cboot}, and is the reason the standard
\SI{0.1}{\micro\farad} bootstrap capacitor must be revisited for this load.

%=====================================================================
\section{Gate-Drive Passive Network Design}\label{sec:passive}

Throughout, the gate charge of one \emph{switch position} is twice that of a single device, because
two MOSFETs are paralleled:
\begin{equation}
Q_{g,\mathrm{pos}} = 2\,Q_{g} = 2\times\SI{211}{\nano\coulomb} = \SI{422}{\nano\coulomb}\ \text{(worst case, to \SI{10}{\volt})}~\cite{ipt015ds}.
\end{equation}
Because the gate is driven to \SI{12}{\volt} rather than the \SI{10}{\volt} at which $Q_g$ is
specified, and because the charge curve is nearly flat above $\sim\SI{7}{\volt}$, a $\approx\SI{10}{\percent}$
margin is added, giving a design value of $Q_{g,\mathrm{pos}}\approx\SI{465}{\nano\coulomb}$ per
switch position.

\subsection{Bootstrap operation and the Miller plateau}\label{sec:boottheory}
\paragraph{Why an N-channel-only half-bridge needs a bootstrap supply}
N-channel MOSFETs are preferred for the half-bridge because, for a given die area and voltage, they
have far lower $\Rdson$ than P-channel devices. But an N-channel MOSFET turns on only when its gate
is several volts above its source. For the low-side device the source is at ground, so a
ground-referenced \SI{12}{\volt} drive works directly. For the high-side device the source is the
switch node, which rises to the full bus voltage when that device is on; to keep it on, the gate
must be driven $\approx\SI{12}{\volt}$ \emph{above the bus} --- a voltage that exists nowhere in a
single-supply system~\cite{ucc27211,balogh}. The bootstrap circuit creates this floating high-side
supply cheaply, using only a diode and a capacitor~\cite{ir_an978}.

\paragraph{How it works}
When the low-side switch is on, the switch node (and the HS pin) is pulled to ground; the bootstrap
capacitor, connected between HB and HS, charges from the \SI{12}{\volt} $V_{DD}$ rail through the
bootstrap diode to $\approx V_{DD}-V_F\approx\SI{11}{\volt}$~\cite{ucc27211,ir_an978}. When the
high-side then turns on, the switch node rises to the bus voltage; the bootstrap diode reverse-biases
(blocking the bus), and the charged capacitor ``rides up'' on the switch node, holding HB
$\approx\SI{11}{\volt}$ above HS and thus supplying the floating bias that keeps the high-side MOSFET
fully enhanced~\cite{ucc27211,ir_an978}. The capacitor is refilled each cycle when the low-side next
turns on --- which is why a minimum low-side on-time (refresh) is required, and why
Section~\ref{sec:cboot} verifies both the droop and the recharge~\cite{ir_an978,ir_an1123}.

\paragraph{The Miller plateau}
During each turn-on the gate voltage rises, pauses on a flat ``Miller plateau''
($\approx\SI{4.4}{\volt}$ here~\cite{ipt015ds}), then rises again. The pause occurs because, once
the channel carries the full load current, all the gate-drive current is diverted into discharging
the gate--drain capacitance $C_{GD}$ so that the drain voltage can swing --- during this interval
$V_{GS}$ is clamped and only the drain voltage moves~\cite{balogh}. The plateau interval is exactly
the window in which the device simultaneously sustains high voltage and high current, so it is where
the switching loss is generated~\cite{balogh}. The Miller charge $Q_{gd}$ that must be pushed through
the plateau is doubled by paralleling ($\approx\SI{102}{\nano\coulomb}$ worst case), and the
\emph{speed} of the voltage transition is set by how much current the driver can force into the gate
at the plateau level --- not by its peak current at full $V_{DD}$~\cite{balogh}. This is why
Balogh emphasizes that the most important attribute of a gate driver is its source/sink capability
\emph{around the plateau}~\cite{balogh}.

\subsection{Bootstrap capacitor}\label{sec:cboot}
The bootstrap capacitor $C_{BOOT}$ must supply the entire charge needed to turn on the high-side
device and to feed the high-side quiescent current during the high-side on-time, while drooping by
no more than the allowed $\Delta V_{BS}$~\cite{ir_an978,ir_an1123}. The governing relations are
\begin{align}
\Delta V_{BS} &\le V_{DD} - V_{F} - V_{GS,\min} - V_{DS,\mathrm{on}}, \label{eq:dvbs}\\
Q_{\mathrm{TOTAL}} &= Q_{g,\mathrm{pos}} + \frac{I_{QBS}+I_{LK}}{\fsw} + Q_{LS},
   \qquad C_{BOOT} \ge \frac{Q_{\mathrm{TOTAL}}}{\Delta V_{BS}}, \label{eq:cboot}
\end{align}
with an additional $\sim\SI{20}{\percent}$ margin recommended for tolerances and for capacitance
loss when the diode is integrated~\cite{ir_an1123}.

\paragraph{Numbers} The boot-diode drop at the relevant charging current is
$V_F\approx\SI{0.9}{\volt}$~\cite{ucc27211}; requiring the high-side gate to stay fully enhanced,
$V_{GS,\min}\approx\SI{10}{\volt}$, with low-side $V_{DS,\mathrm{on}}\approx\SI{0.1}{\volt}$, the
\emph{allowed} droop from \eqref{eq:dvbs} is $\Delta V_{BS}\le 12-0.9-10-0.1=\SI{1.0}{\volt}$. A
conservative design target of $\Delta V_{BS}=\SI{0.5}{\volt}$ is chosen, which keeps the floating
supply at $\approx\SI{10.6}{\volt}$, well above the \SI{6.7}{\volt} high-side UVLO~\cite{ucc27211}.
The quiescent/leakage charge over the worst-case high-side on-time is negligible: with
$I_{QBS}=\SI{0.12}{\milli\ampere}$~\cite{ucc27211}, a high duty $D_{\max}\approx0.95$, and
$t_{HS}=0.95/\SI{40}{\kilo\hertz}=\SI{23.75}{\micro\second}$,
\begin{equation}
Q_{\mathrm{leak}} \approx (\SI{0.12}{\milli\ampere})(\SI{23.75}{\micro\second}) \approx \SI{2.9}{\nano\coulomb} \ll Q_{g,\mathrm{pos}} .
\end{equation}
The gate charge dominates, so $Q_{\mathrm{TOTAL}}\approx\SI{465}{\nano\coulomb}+\SI{3}{\nano\coulomb}\approx\SI{468}{\nano\coulomb}$ and
\begin{equation}
C_{BOOT} \ge \frac{\SI{468}{\nano\coulomb}}{\SI{0.5}{\volt}} \approx \SI{0.94}{\micro\farad}.
\end{equation}
Applying the recommended margin and allowing for DC-bias capacitance loss in Class-II ceramics, the
selected value is $\boxed{C_{BOOT}=\SI{2.2}{\micro\farad},\ \text{X7R},\ \ge\SI{25}{\volt}}$, placed
directly between the HB and HS pins~\cite{ucc27211,ir_an1123}.

\paragraph{Why the datasheet's \SIrange{0.022}{0.1}{\micro\farad} range does not apply here}
That range is correct for a \emph{single} MOSFET with $\sim$30--\SI{90}{\nano\coulomb} of gate
charge~\cite{ucc27211}; the paralleled pair here requires about $5\times$ more charge, which is
exactly why the datasheet states the value ``is dependant on the gate charge of the high-side
MOSFET''~\cite{ucc27211}. Using \SI{0.1}{\micro\farad} would give a droop of
$\SI{468}{\nano\coulomb}/\SI{0.1}{\micro\farad}\approx\SI{4.7}{\volt}$ per cycle --- enough to
under-drive the device or even trip the high-side UVLO. The larger \SI{2.2}{\micro\farad} value is
therefore not optional for this load.

\paragraph{Recharge check} The capacitor replaces only the $\approx\SI{468}{\nano\coulomb}$ removed
each cycle. At the worst duty the low-side on-time is $\approx\SI{1.25}{\micro\second}$; the required
average charging current is $\SI{468}{\nano\coulomb}/\SI{1.25}{\micro\second}\approx\SI{0.37}{\ampere}$,
while the integrated \SI{0.55}{\ohm} boot diode can initially supply
$(12-0.9)/0.55\approx\SI{20}{\ampere}$~\cite{ucc27211}. Recharge is therefore comfortable at
\SI{40}{\kilo\hertz}. At very low motor speed with high-side duty near \SI{100}{\percent}, the
refresh window shrinks and a periodic low-side refresh pulse may be required --- a known limitation
of all bootstrap supplies~\cite{ir_an978,ir_an1123}.

\subsection{Bootstrap diode}
The UCC27211A-Q1 \emph{integrates} the bootstrap diode (anode to $V_{DD}$, cathode to HB); it is
rated \SI{120}{\volt}, has a typical dynamic resistance of \SI{0.55}{\ohm}, a forward drop of
$\approx$0.9--\SI{1.05}{\volt} at \SI{100}{\milli\ampere}, and a \SI{20}{\nano\second} turn-off
time~\cite{ucc27211}. For this design the internal diode is adequate: Section~\ref{sec:cboot} shows
that both the charging current and the droop are comfortably within its capability, and its
\SI{120}{\volt} rating exceeds the \SI{60}{\volt} bus by a wide margin~\cite{ucc27211}. An optional
external diode in parallel may be added to further reduce the forward drop and per-cycle droop; if
used, it must block $\ge\Vbus+V_{DD}$ (a \SI{100}{\volt} part is chosen to match the MOSFET class),
recover quickly ($t_{rr}\le\SI{50}{\nano\second}$) or be a Schottky, and carry the peak charge
current ($\ge\SI{1}{\ampere}$)~\cite{ir_an978,ir_an1123}. Because the function is already on-chip,
this component is an upgrade, not a necessity.

\subsection{Gate resistors and peak-current / Miller check}\label{sec:gateR}
\paragraph{Damping requirement}
The external gate resistor for critically damping the gate-loop resonance formed by the source
inductance $L_S$ and $C_{iss}$ is~\cite{balogh}
\begin{equation}
R_{\mathrm{GATE,OPT}} = 2\sqrt{\frac{L_S}{C_{iss}}} - (R_{DRV}+R_{G,i}).
\end{equation}
For the low-inductance TOLL package ($L_S$ on the order of \SI{1}{\nano\henry}) and
$C_{iss}=\SI{12}{\nano\farad}$, $2\sqrt{L_S/C_{iss}}\approx\SI{0.58}{\ohm}$, which is already
exceeded by the driver output impedance plus the internal gate resistance ($R_{DRV}+R_{G,i}$). The
result is negative, meaning the loop is \emph{already over-damped by the driver impedance and the
internal \SI{1.4}{\ohm} gate resistance alone}~\cite{balogh,ipt015ds}. The external gate resistor is
therefore chosen not for damping but to (i) keep the shared driver output current within rating,
(ii) shape $\mathrm{d}v/\mathrm{d}t$ for EMI, and (iii) balance the two paralleled devices.

\paragraph{Individual resistors}
Best practice for paralleled MOSFETs is an \emph{individual} gate resistor at each gate (two per
switch position) rather than a single shared resistor, so that any tendency of the two devices to
oscillate against one another is damped and their switching is balanced. As a symmetric baseline,
$R_{G,\mathrm{ext}}=\SI{2.2}{\ohm}$ (E24) per MOSFET is close to the \SI{1.8}{\ohm} at which the
IPT015N10N5 switching times are characterized~\cite{ipt015ds}; Section~\ref{sec:asym} refines this
into the asymmetric network that is actually adopted.

\paragraph{Peak-current / $\mathrm{d}v/\mathrm{d}t$ check}
With the gate held at the plateau $V_{\mathrm{plateau}}=\SI{4.4}{\volt}$~\cite{ipt015ds}, the network
would \emph{demand}, per device,
\begin{equation}
I_{G} = \frac{V_{DD}-V_{\mathrm{plateau}}}{R_{G,\mathrm{ext}}+R_{G,i}}
      = \frac{12-4.4}{2.2+1.4} \approx \SI{2.1}{\ampere},
\end{equation}
i.e.\ $\approx\SI{4.2}{\ampere}$ for the pair, which exceeds the \SI{3.7}{\ampere} source limit;
the driver therefore limits the current to \SI{3.7}{\ampere} and itself sets the transition speed.
The time to deliver the combined Miller charge ($2\times\SI{51}{\nano\coulomb}=\SI{102}{\nano\coulomb}$
max) at \SI{3.7}{\ampere} is
\begin{equation}
t_{\mathrm{Miller}} = \frac{Q_{gd,\mathrm{pos}}}{I_{O,\mathrm{source}}}
   = \frac{\SI{102}{\nano\coulomb}}{\SI{3.7}{\ampere}} \approx \SI{28}{\nano\second},
\end{equation}
which, at the worst-case \SI{60}{\volt} bus, gives
\begin{equation}
\frac{\mathrm{d}V_{DS}}{\mathrm{d}t} \approx \frac{\SI{60}{\volt}}{\SI{28}{\nano\second}}
   \approx \SI{2.1}{\volt\per\nano\second},
\end{equation}
which is moderate (good for conducted EMI) and far below the \SI{50}{\volt\per\nano\second} HS
slew-rate limit~\cite{ucc27211}. This confirms the key compatibility question: \emph{even with the
doubled gate charge of paralleling, one channel of the UCC27211A-Q1 switches the pair in under
\SI{30}{\nano\second}} --- precisely the use case for the device's \SI{3.7}{\ampere}/\SI{4.5}{\ampere}
rating~\cite{ucc27211}.

\subsection{Asymmetric (split) gate-resistor network}\label{sec:asym}
A single resistor forces one value to serve both switching edges, which is a compromise. This is
replaced by an \emph{asymmetric} network that sets the turn-on and turn-off speeds independently ---
slower, softer turn-on and fast, low-impedance turn-off --- which is the configuration adopted in
the final design.

\paragraph{Topology}
Per MOSFET, an always-present series resistor $R_{on}$ carries the turn-on (source) current, and a
branch of $R_{off}$ in series with an OFF Schottky diode $D_{off}$ is placed in parallel with it.
$D_{off}$ is oriented anode-to-gate, cathode-to-driver-output, so it conducts \emph{only when the
gate is pulled down} (turn-off / $\mathrm{d}v/\mathrm{d}t$ clamp). Hence
\begin{equation}
R_{g(\mathrm{on})} = R_{on}, \qquad R_{g(\mathrm{off})} = R_{on}\parallel(R_{off}+R_{D}),
\end{equation}
so at turn-on the gate sees only $R_{on}$, and at turn-off it sees the lower parallel
combination~\cite{balogh,pmeg4030er}.

\paragraph{Design direction --- slow on, fast off}
For a hard-switched half-bridge the preferred asymmetry is slower turn-on, faster
turn-off~\cite{pmeg4030er,an90059,toshiba_dvdt}. Faster turn-on raises the $\mathrm{d}i/\mathrm{d}t$
of the drain current, which increases the reverse-recovery current of the complementary device's
freewheeling body diode; the IPT015N10N5 body diode stores $Q_{rr}\approx\SI{316}{\nano\coulomb}$,
so deliberately softened turn-on limits that recovery peak, its associated loss, and the EMI it
generates~\cite{ipt015ds,pmeg4030er,an90059}. Turn-off, conversely, benefits from being fast and
low-impedance: it reduces turn-off loss and, crucially, holds the gate of the \emph{off} device
firmly at ground so that the Miller current injected through $C_{GD}$ when the complementary switch
slews the node cannot raise $V_{GS}$ to the threshold and cause self (false)
turn-on~\cite{an90059,toshiba_dvdt}. The \SI{100}{\volt} device on a bus up to \SI{60}{\volt} has
enough voltage reserve ($\approx\SI{40}{\volt}$) to absorb the slightly higher turn-off overshoot
that fast turn-off produces; because that reserve is smaller than on a fixed \SI{48}{\volt} bus, the
ability to control the turn-off edge independently is all the more valuable. The asymmetry also
reinforces the driver's own \SI{4.5}{\ampere}-sink / \SI{3.7}{\ampere}-source
imbalance~\cite{ucc27211}.

\paragraph{Value selection}
\begin{itemize}[leftmargin=1.4em]
  \item $R_{on}=\SI{3.3}{\ohm}$ (per MOSFET). Raised above the \SI{2.2}{\ohm} baseline to soften the
        turn-on edge. The resulting per-device plateau current is
        $(12-4.4)/(R_{on}+R_{G,i})=7.6/(3.3+1.4)=\SI{1.62}{\ampere}$ ($\approx\SI{3.2}{\ampere}$ for
        the pair) --- now set by the resistor rather than the \SI{3.7}{\ampere} driver limit --- giving
        a controlled $\mathrm{d}v/\mathrm{d}t$ of $60/(102\,\mathrm{n}/3.2)\approx\SI{1.9}{\volt\per\nano\second}$
        worst case~\cite{ipt015ds,ucc27211}.
  \item $R_{off}=\SI{1.0}{\ohm}$ (per MOSFET). With the OFF-diode dynamic resistance
        $R_{D}\approx$0.1--\SI{0.15}{\ohm} (from the PMEG4030ER conduction characteristic,
        $V_F\approx$0.46--\SI{0.54}{\volt})~\cite{pmeg4030er},
        $R_{g(\mathrm{off})}=R_{on}\parallel(R_{off}+R_D)=3.3\parallel1.15\approx\SI{0.85}{\ohm}$, an
        effective on:off resistance ratio of about $3.9\!:\!1$.
\end{itemize}
The internal gate resistance $R_{G,i}=\SI{1.4}{\ohm}$ is in series with the die and cannot be
bypassed~\cite{ipt015ds}, so the total turn-off gate-loop resistance is
$R_{G,i}+R_{g(\mathrm{off})}=\SI{2.25}{\ohm}$; reducing $R_{off}$ below $\sim\SI{1}{\ohm}$ yields
diminishing returns, so \SI{1.0}{\ohm} is well placed. At this value the per-device turn-off plateau
current is $4.4/2.25\approx\SI{1.96}{\ampere}$ ($\approx\SI{3.9}{\ampere}$ for the pair, near the
\SI{4.5}{\ampere} sink limit), giving a turn-off transition of $\approx\SI{26}{\nano\second}$ versus
$\approx\SI{32}{\nano\second}$ at turn-on --- the desired asymmetry~\cite{ucc27211}.

\paragraph{OFF diode selection}
$D_{off}$ must block the reverse voltage that appears across $R_{on}$ at turn-on
($\approx I_{on,\mathrm{pk}}R_{on}\approx 1.62\times3.3\approx\SI{5.3}{\volt}$ plus ringing), carry
the turn-off current pulse ($\approx\SI{2}{\ampere}$ peak per gate), recover instantly, and have low
dynamic resistance. A Schottky satisfies all four. The selected component is the \textbf{Nexperia
PMEG4030ER} --- \SI{40}{\volt}, \SI{3}{\ampere} Schottky rectifier with low $V_F$ ($\approx$0.46--\SI{0.54}{\volt}),
$I_{FSM}=\SI{50}{\ampere}$, SOD-123W~\cite{pmeg4030er}. Its \SI{40}{\volt} rating gives $>7\times$
margin over the $\sim\SI{5}{\volt}$ reverse stress, and its low $V_F$ and small junction capacitance
keep the turn-off path clean~\cite{pmeg4030er}.

\paragraph{Resistor power and package}
The average power in a gate resistor is a fraction of the total gate-drive loss
$P_G=Q_g V_{DD}\fsw$~\cite{balogh}. Per device, with $Q_g\approx\SI{232}{\nano\coulomb}$ at
\SI{12}{\volt}, $P_G\approx\SI{0.11}{\watt}$, split between turn-on and turn-off; the external
resistor sees only its fraction of the gate-loop resistance, so $P_{R_{on}}\approx\SI{26}{\milli\watt}$
and $P_{R_{off}}\approx\SI{6}{\milli\watt}$ --- both far below 1/16\,W. The resistor does see
short current pulses ($\sim$1.6--\SI{2}{\ampere}), but the per-transition energy in $R_{on}$ is
$\approx\SI{0.6}{\micro\joule}$, negligible against any chip resistor's pulse rating. A
\textbf{0805 thick-film, 1/8\,W} package is chosen for both $R_{on}$ and $R_{off}$
($>4\times$ average-power margin and good pulse robustness)~\cite{balogh}.

\subsection{Gate--source pull-down resistors}\label{sec:pulldown}
A gate-to-source resistor holds each MOSFET definitely off before the driver outputs become active
at power-up, and adds $\mathrm{d}v/\mathrm{d}t$ protection against Miller-induced
turn-on~\cite{balogh}. In normal operation the driver's low output impedance and \SI{4.5}{\ampere}
sink dominate this protection; the pull-down's main job is the power-up state. A standard
$\boxed{R_{GS}=\SI{10}{\kilo\ohm}}$ per MOSFET is large enough not to load the driver and small
enough to hold the gate down at power-up. A capacitive-divider estimate confirms the margin: when
the complementary switch imposes the worst-case \SI{60}{\volt} on the off device, the worst-case
gate rise through the divider $C_{rss}/C_{iss}=\SI{80}{\pico\farad}/\SI{12}{\nano\farad}\approx0.0067$
would be only $\approx\SI{0.40}{\volt}$ even with a \emph{floating} gate --- still far below the
\SI{2.2}{\volt} threshold --- and the active \SI{4.5}{\ampere} sink holds the real value lower
still~\cite{ipt015ds,ucc27211}. A \textbf{0603, 1/10\,W} package suffices (the resistor
dissipates $12^2/\SI{10}{\kilo\ohm}=\SI{14.4}{\milli\watt}$ when the gate is high).

\subsection{Decoupling capacitors}
\paragraph{$V_{DD}$ (low-side) decoupling}
Whenever the LO output sources gate current, an equal pulse is drawn from $V_{DD}$, so a low-ESR
ceramic must sit directly between $V_{DD}$ and $V_{SS}$~\cite{ucc27211}. Sized for a tolerable ripple
$\Delta V$ with the same charge argument as the bootstrap,
$C_{VDD}\ge Q_{g,\mathrm{pos}}/\Delta V=\SI{465}{\nano\coulomb}/\SI{0.5}{\volt}\approx\SI{0.93}{\micro\farad}$.
TI recommends \SIrange{0.22}{4.7}{\micro\farad}~\cite{ucc27211}; the chosen value is
$\boxed{C_{VDD}=\SI{4.7}{\micro\farad}\ \text{X7R}+\SI{0.1}{\micro\farad}}$ per driver, one set for
each of the three drivers.

\paragraph{HB (high-side) decoupling}
This role is performed by the bootstrap capacitor itself (Section~\ref{sec:cboot}), placed directly
between HB and HS; the separate \SIrange{0.022}{0.1}{\micro\farad} recommendation is only the minimum
HF decoupling and is exceeded here by the computed \SI{2.2}{\micro\farad}~\cite{ucc27211}.

\paragraph{\SI{12}{\volt} bulk}
A shared bulk reservoir (e.g.\ \SIrange{22}{47}{\micro\farad}) on the auxiliary \SI{12}{\volt} rail
feeds the three drivers and absorbs the aggregate gate-drive pulses~\cite{ucc27211}.

\paragraph{DC-link decoupling}
Because the source is a battery connected through a cable harness, the source impedance (internal
battery resistance plus harness inductance) is significant, so local DC-link decoupling is even more
important than on a stiff laboratory bus. The power stage carries a ripple current on the order of
\SI{100}{\ampere} (up to $\approx\SI{132}{\ampere}$ worst case at the depleted \SI{24}{\volt} pack)
and must be decoupled by a low-ESR/ESL bank (film and/or ceramic) close to the half-bridge, sized
for the RMS ripple current and to keep the commutation overshoot within the \SI{100}{\volt} rating.
Since the highest battery voltage is \SI{60}{\volt}, the bus capacitors should be rated
$\ge\SI{80}{\volt}$ (film) or $\ge\SI{100}{\volt}$ (ceramic, allowing for DC-bias derating). Precise
sizing depends on the harness inductance and is a board-level task.

\begin{table}[!t]
\centering
\caption{Gate-drive passive bill of materials (per leg / per board).}
\label{tab:gdbom}
\renewcommand{\arraystretch}{1.2}
\begin{tabularx}{\textwidth}{@{}l L L c c@{}}
\toprule
\textbf{Ref.} & \textbf{Function} & \textbf{Value} & \textbf{Per leg} & \textbf{3 legs} \\
\midrule
$C_{BOOT}$ & Bootstrap capacitor (HB--HS) & \SI{2.2}{\micro\farad} X7R $\ge\SI{25}{\volt}$ & 1 & 3 \\
$C_{VDD}$ & $V_{DD}$ decoupling & \SI{4.7}{\micro\farad} X7R $+$ \SI{0.1}{\micro\farad} & 1 set & 3 sets \\
$R_{on}$ & Turn-on resistor (per MOSFET) & \SI{3.3}{\ohm}, 0805, 1/8\,W & 4 & 12 \\
$R_{off}$ & Turn-off resistor (per MOSFET) & \SI{1.0}{\ohm}, 0805, 1/8\,W & 4 & 12 \\
$D_{off}$ & OFF Schottky diode & Nexperia PMEG4030ER (\SI{40}{\volt}, \SI{3}{\ampere}) & 4 & 12 \\
$R_{GS}$ & Gate--source pull-down & \SI{10}{\kilo\ohm}, 0603, 1/10\,W & 4 & 12 \\
$D_{BOOT}$ (opt.) & External parallel boot diode & \SI{100}{\volt} fast/Schottky $\ge\SI{1}{\ampere}$ & 0--1 & 0--3 \\
\bottomrule
\end{tabularx}
\end{table}

\noindent Plus the active devices: \textbf{$12\times$ IPT015N10N5} (four per leg) and
\textbf{$3\times$ UCC27211A-Q1}.

%=====================================================================
\section{Phase-Current Sensing}\label{sec:current}

\subsection{Survey and comparison of current-sense topologies}\label{sec:csurvey}
Four families of inverter current sensing exist; their measurement node is shown in
Fig.~\ref{fig:topo} for one half-bridge, and they are compared in Table~\ref{tab:topo}.

\begin{figure}[!t]
\centering
\begin{circuitikz}[scale=0.92,transform shape,american]
\draw (0,4) node[anchor=east]{$\Vbus$} -- (6.5,4);
\draw (0,-1.2) node[anchor=east]{GND} -- (6.5,-1.2);
\draw (2,3.0) rectangle (3,3.8) node[midway]{$Q_H$};
\draw (2.5,4) -- (2.5,3.8);
\draw (2.5,3.0) -- (2.5,1.7);
\draw (2,0.9) rectangle (3,1.7) node[midway]{$Q_L$};
\draw (2.5,1.7) node[circ]{} -- (4.6,1.7) node[anchor=west]{phase $\to$ motor};
\draw (3.6,1.7) node[anchor=south]{\scriptsize in-line $\Rshunt$};
\draw[fill=yellow!30] (3.25,1.55) rectangle (3.95,1.85);
\draw (2.5,0.9) -- (2.5,0.2);
\draw[fill=yellow!30] (2.15,-0.1) rectangle (2.85,0.2) node[midway]{};
\draw (2.5,0.0) node[anchor=west]{\scriptsize low-side $\Rshunt$};
\draw (2.5,-0.1) -- (2.5,-1.2);
\draw[fill=yellow!30] (2.15,4.0) rectangle (2.85,4.3);
\draw (2.5,4.45) node[anchor=south]{\scriptsize high-side $\Rshunt$ (bus)};
\draw (2.5,4.3) -- (2.5,5.0) -- (5.5,5.0) -- (5.5,4);
\draw[fill=orange!25] (5.15,-1.55) rectangle (5.85,-1.25);
\draw (5.5,-1.7) node[anchor=north]{\scriptsize single DC-link $\Rshunt$};
\draw (5.5,-1.2)--(5.5,-1.25); \draw(5.5,-1.55)--(5.5,-2.0)--(0,-2.0);
\end{circuitikz}
\caption{Possible shunt positions in one inverter leg: high-side (in series with the leg from the
bus), low-side (in series with the leg to ground), in-line (in series with the phase conductor),
and the single DC-link shunt in the common return. The Hall/magnetic alternative replaces the shunt
with a field sensor around the phase conductor.}
\label{fig:topo}
\end{figure}

\begin{table}[!t]
\centering
\caption{Comparison of current-sense topologies for a three-phase inverter.}
\label{tab:topo}
\renewcommand{\arraystretch}{1.25}
\begin{tabularx}{\textwidth}{@{}p{2.5cm} L L L L@{}}
\toprule
\textbf{Attribute} & \textbf{Low-side shunt} & \textbf{High-side shunt} & \textbf{In-line (in-phase) shunt} & \textbf{Magnetic / Hall} \\
\midrule
Common-mode voltage at sense element & Near ground ($\approx0$) & At/near $\Vbus$ & Swings rail-to-rail at $\fsw$ with high $\mathrm{d}v/\mathrm{d}t$ & Galvanically isolated; no CM \\
Measures true phase current at all duty cycles & No --- only while $Q_L$ conducts; blind near \SI{100}{\percent} duty & No --- only while $Q_H$ conducts & \textbf{Yes --- continuous, at any duty} & Yes \\
Amplifier requirement & Low-CM CSA & High-CM CSA & High-CM CSA with PWM rejection & --- (sensor is the amplifier) \\
DC capability & Yes & Yes & Yes & Yes (Hall) / No (CT, Rogowski) \\
Suitability for FOC & Limited (reconstruction needed) & Limited & \textbf{Direct (per-phase, continuous)} & Good but costlier/larger \\
\bottomrule
\end{tabularx}
\end{table}

For field-oriented control the inner current loop needs the \emph{actual} phase current
continuously and independently of duty cycle. The low-side and high-side shunts measure current
only while their switch conducts, becoming blind near \SI{100}{\percent} and \SI{0}{\percent} duty
respectively; the single DC-link shunt requires duty-dependent reconstruction. The \emph{in-line}
(in-phase) shunt, placed in series with the phase conductor, observes the true phase current at any
duty cycle, which is why it is chosen here, in line with TI's TIDA-00913 reference
design~\cite{tida00913}. Its cost is a sense element whose common-mode potential swings from
rail to rail at the switching frequency with high $\mathrm{d}v/\mathrm{d}t$, which dictates the
amplifier requirements in Section~\ref{sec:csa}.

\subsection{Shunt selection}\label{sec:shunt}
A metal-element (metal-strip) shunt is used for its low inductance and low temperature coefficient.
The shunt voltage is $v_{\mathrm{sh}}(t)=\Rshunt\,i(t)+\mathrm{ESL}\,\mathrm{d}i/\mathrm{d}t$; at
\SI{40}{\kilo\hertz} with fast switching edges the inductive term must be negligible against the
\SI{0.1}{\milli\ohm} resistive term, which a metal-strip element with $<\SI{2}{\nano\henry}$
achieves. At the worst-case continuous current ($\approx\SI{119}{\ampere}$ RMS) the shunt dissipates
$I^2\Rshunt\approx119^2\times\SI{0.1}{\milli\ohm}\approx\SI{1.4}{\watt}$; a \SI{0.1}{\milli\ohm},
\SI{5}{\watt} element gives ample margin. An alternative is two extremely common
\SI{0.2}{\milli\ohm}, 2512 elements in parallel ($=\SI{0.1}{\milli\ohm}$, each dissipating
$\approx\SI{0.7}{\watt}$).

\subsection{Current-sense amplifier and gain}\label{sec:csa}
The amplifier must (i) tolerate a common-mode voltage that swings from slightly below ground
(body-diode conduction during dead-time, $\approx-\SI{1}{\volt}$) to $\Vbus=\SI{60}{\volt}$ plus
ringing overshoot; (ii) reject the large common-mode $\mathrm{d}v/\mathrm{d}t$ at each switching
edge so the output does not ring; (iii) be \emph{bidirectional}, since the phase current is AC; and
(iv) have a bandwidth well above \SI{40}{\kilo\hertz}. The Texas Instruments INA240 satisfies all
four: it is a voltage-output, zero-drift current-sense amplifier with a $-4$ to $+\SI{80}{\volt}$
common-mode range \emph{independent of supply}, with enhanced PWM rejection explicitly designed to
suppress the common-mode PWM transients of motor drives, offered in four fixed gains
(\num{20}, \num{50}, \num{100}, \num{200}\,V/V)~\cite{ina240ds}. The $-\SI{4}{\volt}$ lower limit
accepts the sub-ground excursion, and the $+\SI{80}{\volt}$ upper limit leaves \SI{20}{\volt} of
reserve above the \SI{60}{\volt} bus for ringing. The PWM-rejection benefit is application-dependent
rather than a single specified frequency; TI's validated TIDA-00913 reference design demonstrates
accurate in-line operation up to PWM switching frequencies of \SI{100}{\kilo\hertz}, so the
\SI{40}{\kilo\hertz} SVPWM here is comfortably within the validated range~\cite{ina240ds,tida00913}.

On a \SI{3.3}{\volt} analog supply, the bidirectional output at zero current is centered by the two
reference pins of the INA240, whose internal divider sets the output midpoint at
$V_{\mathrm{ref}}=(V_{\mathrm{REF1}}+V_{\mathrm{REF2}})/2$~\cite{ina240ds}; tying REF1 to ground and
REF2 to the \SI{3.3}{\volt} rail gives $V_{\mathrm{ref}}=\SI{1.65}{\volt}$. The output is then
\begin{equation}
V_{\mathrm{out}} = V_{\mathrm{ref}} + G\,\Rshunt\,i_{\mathrm{ph}} .
\label{eq:csaout}
\end{equation}
With a rail-to-rail output the usable single-sided swing is $\approx\SI{1.45}{\volt}$ (\SI{0.2}{\volt}
margin to each rail). Choosing gain $G=50$ (INA240A2), the linear full-scale peak current is
\begin{equation}
\hat I_{\mathrm{FS}} = \frac{1.45}{G\,\Rshunt}
   = \frac{1.45}{50\times\SI{0.1}{\milli\ohm}} = \SI{290}{\ampere}\ \text{peak}\ (\SI{205}{\ampere}\,\text{RMS}).
\end{equation}
The worst-case peak of \SI{169}{\ampere} produces an output excursion of
$G\Rshunt\hat I_{\mathrm{ph}}=50\times\SI{0.1}{\milli\ohm}\times169=\SI{0.85}{\volt}$, i.e.\
$V_{\mathrm{out}}\in[0.80,2.50]\,\si{\volt}$, leaving $\approx\SI{0.5}{\volt}$ before clipping. This
reserve is deliberately kept for transient overshoot and lets an over-current (OC) threshold be set
at $\approx\SI{220}{\ampere}$ peak ($\approx1.3\times\hat I_{\mathrm{ph}}^{\max}$), detectable by an
STM32G4 analog comparator or in software before the amplifier saturates. The quantization and offset,
referred to current, are
\begin{align}
\Delta I_{\mathrm{LSB}} &= \frac{V_{\mathrm{REF+}}/2^{N}}{G\,\Rshunt}
   = \frac{3.3/4096}{50\times\SI{0.1}{\milli\ohm}} = \SI{0.16}{\ampere/LSB},\\
I_{\mathrm{os}} &= \frac{V_{\mathrm{os(max)}}}{\Rshunt} = \frac{\SI{25}{\micro\volt}}{\SI{0.1}{\milli\ohm}}
   = \SI{0.25}{\ampere},
\end{align}
both negligible against a \SI{119}{\ampere} signal (the INA240 zero-drift architecture gives
$V_{\mathrm{os}}\le\SI{25}{\micro\volt}$)~\cite{ina240ds}. The gain-50 / \SI{0.1}{\milli\ohm} pair is
therefore preferred over a gain-20 / \SI{0.25}{\milli\ohm} pair, which would dissipate
$\approx\SI{4.6}{\watt}$ per shunt for the same full scale.

\subsection{ADC interface and output filtering}\label{sec:csafilter}
For FOC the phase current is sampled \emph{synchronously} with the PWM carrier (typically at the
counter peak/valley) so that the sampling instant coincides with the middle of the switching state,
where the inductor ripple is at its average value; this naturally rejects the \SI{40}{\kilo\hertz}
ripple and its harmonics without a heavy filter~\cite{an5346}. A heavy anti-alias filter would
instead add phase lag to the very signal on which the current loop closes, degrading the loop. The
CSA output therefore uses only a light first-order RC filter ($R_f=\SI{100}{\ohm}$,
$C_f=\SI{2.2}{\nano\farad}$) that (i) band-limits the high-frequency switching noise and ringing
(corner $f_c=1/2\pi R_fC_f=\SI{723}{\kilo\hertz}$, far above the kHz-range loop content but damping
the tens-of-MHz ringing) and (ii) acts as a charge reservoir for the ADC sample-and-hold. The low
output impedance of the INA240 drives this network and the ADC sampling capacitor with ample settling
margin. STM32G4 hardware oversampling/averaging may be enabled for further noise reduction at the cost
of bandwidth~\cite{an5346}. Fig.~\ref{fig:csa} shows the complete per-phase chain.

\begin{figure}[!t]
\centering
\begin{circuitikz}[scale=1.0,transform shape,american]
\draw (0,4.7) node[anchor=east]{phase node} -- (1.3,4.7);
\draw[fill=yellow!30] (1.3,4.5) rectangle (2.5,4.9);
\draw (1.9,5.2) node{\scriptsize $\Rshunt=\SI{0.1}{\milli\ohm}$ (Kelvin)};
\draw (2.5,4.7) -- (3.9,4.7) node[anchor=west]{to motor};
\draw (2.3,4.5) -- (2.3,3.0) -- (2.7,3.0);
\draw (1.5,4.5) -- (1.5,2.5) -- (2.7,2.5);
\draw (2.45,3.0) node[anchor=south]{\scriptsize IN+};
\draw (2.0,2.5) node[anchor=south]{\scriptsize IN--};
\draw (2.7,1.9) rectangle (5.1,3.5);
\draw (3.9,2.45) node{\small INA240A2};
\draw (3.9,3.05) node{\scriptsize $G=50$};
\draw (3.4,3.5) -- (3.4,4.0) node[anchor=east]{\scriptsize $3V3_{\mathrm A}$+\SI{100}{\nano\farad}};
\draw (3.9,1.9) -- (3.9,1.2) node[anchor=north,align=center]{\scriptsize REF1=GND, REF2=3V3\\\scriptsize $\Rightarrow V_{\mathrm{ref}}=\SI{1.65}{\volt}$};
\draw (5.1,2.7) -- (5.6,2.7) to[R,l={$R_f$, \SI{100}{\ohm}}] (7.1,2.7);
\draw (7.1,2.7) node[circ]{} to[C,l={$C_f$, \SI{2.2}{\nano\farad}}] (7.1,1.5) node[ground]{};
\draw (7.1,2.7) -- (7.8,2.7) node[anchor=west]{\scriptsize ADC};
\end{circuitikz}
\caption{In-line per-phase current-sense chain: Kelvin-connected \SI{0.1}{\milli\ohm} metal-strip
shunt, INA240A2 ($G=50$) with mid-scale reference, and a light \SI{723}{\kilo\hertz} RC into the
STM32G4 ADC. Three identical chains are used (phases U, V, W).}
\label{fig:csa}
\end{figure}

%=====================================================================
\section{Voltage Sensing}\label{sec:voltage}

\subsection{Requirements}
Hardware voltage measurement is needed for (i) the DC-bus voltage --- for over-/under-voltage
protection, modulation-index scaling, and power calculation --- and (ii) each of the three phase-node
voltages --- for output monitoring, phase-loss / motor-disconnect detection, and diagnostics. The
phase-voltage chains also leave a hardware path open for future back-EMF-based position observation
without a board change, although the present design commutates from Hall sensors and does not rely
on phase-voltage feedback for control. All four channels share the inverter ground (the design is
non-isolated and uses a single shared ground), so each is a grounded resistive divider feeding an
STM32G4 ADC channel. The bus channel is slow (protection/scaling); the phase channels preserve more
bandwidth for monitoring and for the optional back-EMF capability.

\subsection{Resistive-divider design}\label{sec:divider}
The ADC conversion range is $0$ to $V_{\mathrm{REF+}}$; using
$V_{\mathrm{REF+}}=V_{\mathrm{DDA}}=\SI{3.3}{\volt}$ maximizes the range. The divider ratio
$k=(R_{\mathrm{top}}+R_{\mathrm{bot}})/R_{\mathrm{bot}}$ maps the maximum measured voltage to near
full scale with reserve. Choosing $R_{\mathrm{top}}=\SI{110}{\kilo\ohm}$ and
$R_{\mathrm{bot}}=\SI{4.7}{\kilo\ohm}$,
\begin{equation}
k=\frac{R_{\mathrm{top}}+R_{\mathrm{bot}}}{R_{\mathrm{bot}}}=\frac{110+4.7}{4.7}=24.4,
\qquad V_{\mathrm{ADC}}=\frac{V_{\mathrm{in}}}{k}.
\end{equation}
This maps $\Vbus=\SI{60}{\volt}\to\SI{2.46}{\volt}$ and $\SI{30}{\volt}\to\SI{1.23}{\volt}$, with
full scale at $\approx\SI{73}{\volt}$ --- enough headroom to read transients above the \SI{60}{\volt}
nominal maximum. The same ratio is used for the phase dividers, since a phase node can reach the bus
voltage. The standing current and dissipation at \SI{60}{\volt} are
\begin{equation}
I_{\mathrm{div}}=\frac{60}{R_{\mathrm{top}}+R_{\mathrm{bot}}}=\SI{523}{\micro\ampere},\quad
P_{R_{\mathrm{top}}}=I_{\mathrm{div}}^2R_{\mathrm{top}}=\SI{30}{\milli\watt},
\end{equation}
so all four dividers together draw $\approx\SI{126}{\milli\watt}$ --- negligible. The top resistor
sees the full bus voltage; a single thick-film \SI{110}{\kilo\ohm} chip rated $\ge\SI{150}{\volt}$
(e.g.\ 0805/1206) is adequate at \SI{60}{\volt}, but splitting $R_{\mathrm{top}}$ into two series
\SI{55}{\kilo\ohm} parts is recommended to halve the per-component voltage stress and improve ratio
matching. The Th\'evenin source impedance presented to the ADC is
\begin{equation}
R_{\mathrm{src}}=R_{\mathrm{top}}\parallel R_{\mathrm{bot}}
   =\frac{110\text{k}\times4.7\text{k}}{110\text{k}+4.7\text{k}}=\SI{4.51}{\kilo\ohm},
\label{eq:rsrc}
\end{equation}
which is the key quantity for the buffer question of Section~\ref{sec:buffer}.

\subsection{Does the ADC need an op-amp buffer? A quantitative answer}\label{sec:buffer}
An STM32G4 ADC channel samples by connecting the source, through an internal multiplexer/switch
resistance $R_{\mathrm{ADC}}$, to a sampling capacitor $C_{\mathrm{ADC}}$ for the programmed sampling
time $t_s$ (Fig.~\ref{fig:adcmodel}). The capacitor must settle to within $\tfrac12$\,LSB of the
source voltage. For an $N$-bit conversion this requires~\cite{an2834,an5346}
\begin{equation}
t_s \;\ge\; \bigl(R_{\mathrm{AIN}}+R_{\mathrm{ADC}}\bigr)\,C_{\mathrm{ADC}}\,\ln\!\bigl(2^{\,N+1}\bigr),
\label{eq:settle}
\end{equation}
where $R_{\mathrm{AIN}}=R_{\mathrm{src}}$ is the external source impedance. ST does not publish
$R_{\mathrm{ADC}}$/$C_{\mathrm{ADC}}$ directly for the G4; instead the datasheet gives a ``maximum
$R_{\mathrm{AIN}}$ versus sampling time'' table and recommends choosing the sampling time from it for
the given source impedance~\cite{an5346}. Using the AN2834 settling model with representative values
$C_{\mathrm{ADC}}\approx\SI{5}{\pico\farad}$, $R_{\mathrm{ADC}}\approx\SI{2}{\kilo\ohm}$, and
$\ln(2^{13})=9.01$ for the 12-bit mode,
\begin{equation}
t_s \ge (4510+2000)(5\times10^{-12})(9.01)=\SI{293}{\nano\second}.
\end{equation}
At a \SI{60}{\mega\hertz} ADC clock this is \num{17.6} cycles, whereas the STM32G4 offers
sampling-time settings up to \num{640.5} cycles; the \num{47.5}-cycle setting alone is
\SI{792}{\nano\second} --- more than $2.7\times$ the requirement. Because the bus and phase voltages
need only kHz-rate refresh, spending a longer sampling time costs nothing. Therefore:
\begin{center}
\emph{No external op-amp buffer is required for the voltage dividers.}
\end{center}
The divider source impedance settles to 12-bit accuracy well within a trivially available sampling
time; a small filter capacitor (Section~\ref{sec:vfilter}) further reduces the effective source
impedance at the sampling instant.

\begin{figure}[!t]
\centering
\begin{circuitikz}[scale=1.0,transform shape,american]
\draw (0,1.2) node[anchor=east]{$V_{\mathrm{div}}$} to[R,l={$R_{\mathrm{AIN}}{=}R_{\mathrm{src}}$}] (3,1.2);
\draw (3,1.2) to[R,l=$R_{\mathrm{ADC}}$] (5,1.2);
\draw (5,1.2) to[nos,l=$t_s$] (6.2,1.2);
\draw (6.2,1.2) to[C,l=$C_{\mathrm{ADC}}$] (6.2,0) node[ground]{};
\draw (6.2,1.2) -- (7.0,1.2) node[anchor=west]{SAR};
\end{circuitikz}
\caption{First-order model of the STM32G4 ADC input stage used in \eqref{eq:settle}: the divider
source impedance $R_{\mathrm{src}}$ in series with the internal switch resistance charges the
sampling capacitor $C_{\mathrm{ADC}}$ over the sampling window $t_s$.}
\label{fig:adcmodel}
\end{figure}

\begin{table}[!t]
\centering
\caption{Op-amp buffer before the ADC: benefits versus costs.}
\label{tab:buffer}
\renewcommand{\arraystretch}{1.2}
\begin{tabularx}{\textwidth}{@{}L L@{}}
\toprule
\textbf{Arguments for a buffer} & \textbf{Arguments against a buffer}\\
\midrule
Low output impedance $\Rightarrow$ fast settling, short sampling time, high multiplexed scan rate &
Adds its own offset, drift, noise, and finite-bandwidth error to an \emph{absolute} measurement\\
Isolates the divider from ADC charge-injection ``kickback'' and inter-channel crosstalk &
Extra BOM cost, board area, and supply decoupling\\
Enables an active anti-alias filter &
Becomes an additional calibration burden\\
Prevents resistive loading of the divider &
Unnecessary here: the required $t_s$ is small, \eqref{eq:settle}\\
\bottomrule
\end{tabularx}
\end{table}

If a buffer were ever desired --- for example to share one divider among several fast multiplexed
channels, or to add an active filter --- it should be realized with the \emph{MCU's own resources}
rather than an external IC: the STM32G473 contains operational amplifiers that can be configured as
unity-gain followers (or PGA) with their output routed \emph{internally} to an ADC channel, which
both removes the external-pin noise-pickup path and costs nothing on the BOM~\cite{an5306,stm32g473ds}.
ST documents exactly this follower-to-ADC use and recommends a \SI{200}{\nano\second} sampling time
when reading an internal op-amp output~\cite{an5346,an5306}. For the present low-rate voltage
channels the conclusion is to omit the buffer and use a longer sampling time plus a filter capacitor.

\subsection{Bus/phase filtering and protection}\label{sec:vfilter}
A capacitor $C_{\mathrm{f}}$ in parallel with $R_{\mathrm{bot}}$ forms a first-order low-pass filter
with corner $f_c=1/(2\pi R_{\mathrm{src}}C_{\mathrm{f}})$.
\textbf{Bus channel:} $C_{\mathrm{f}}=\SI{22}{\nano\farad}$ gives $f_c=\SI{1.6}{\kilo\hertz}$, removing
the switching ripple from the slow bus measurement. \textbf{Phase channels:} the filter must reject
the \SI{40}{\kilo\hertz} PWM while passing the fundamental of the monitored phase voltage; a corner
of $f_c\approx$\,2--\SI{3}{\kilo\hertz} attenuates the carrier at the cost of a phase lag
$\phi=\arctan(f_e/f_c)$, where $f_e$ is the electrical frequency. If the phase voltages are used for
back-EMF observation (the provisioned future capability), this lag is deterministic and is
compensated in firmware (the observer shifts the sampled phase by $\phi$); Table~\ref{tab:phase}
lists representative values. \textbf{Protection:} during dead-time/body-diode conduction a phase node
swings to $\approx-\SI{1}{\volt}$, scaled by the divider to $-1/24.4=-\SI{41}{\milli\volt}$ at the
ADC pin --- small, but a Schottky clamp (e.g.\ BAT54S) from the ADC node to ground and to
$V_{\mathrm{DDA}}$ is added as insurance against over-/under-range transients.

\begin{table}[!t]
\centering
\caption{Phase low-pass filter delay $\phi=\arctan(f_e/f_c)$ (compensated in firmware if back-EMF observation is used).}
\label{tab:phase}
\renewcommand{\arraystretch}{1.2}
\begin{tabular}{c c c c}
\toprule
$f_c$ & $f_e=\SI{200}{\hertz}$ & $f_e=\SI{500}{\hertz}$ & $f_e=\SI{1000}{\hertz}$\\
\midrule
\SI{2}{\kilo\hertz} & \ang{5.7} & \ang{14.0} & \ang{26.6}\\
\SI{3}{\kilo\hertz} & \ang{3.8} & \ang{9.5}  & \ang{18.4}\\
\bottomrule
\end{tabular}
\end{table}

\begin{figure}[!t]
\centering
\begin{circuitikz}[scale=1.0,transform shape,american]
\draw (0,3.2) node[anchor=east]{$V_{\mathrm{in}}$ (bus/phase)} -- (1.0,3.2);
\draw (1.0,3.2) to[R,l=$R_{\mathrm{top}}$,a=\SI{110}{\kilo\ohm}] (1.0,1.6);
\draw (1.0,1.6) node[circ]{} -- (2.6,1.6) node[anchor=west]{to ADC};
\draw (1.0,1.6) to[R,l=$R_{\mathrm{bot}}$,a=\SI{4.7}{\kilo\ohm}] (1.0,0.0) node[ground]{};
\draw (2.0,1.6) to[C,l=$C_f$] (2.0,0.0) node[ground]{};
\draw (2.6,1.6) -- (3.2,1.6);
\draw (3.2,1.6) to[Do] (3.2,2.8) node[anchor=south]{\scriptsize $V_{\mathrm{DDA}}$};
\draw (3.2,1.6) node[circ]{}; \draw (3.6,1.6) to[Do] (3.6,0.4) node[ground]{};
\draw (3.4,2.2) node[anchor=west]{\scriptsize BAT54S};
\end{circuitikz}
\caption{Measurement divider (ratio \num{24.4}), first-order filter capacitor, and Schottky clamp to
$V_{\mathrm{DDA}}$/GND. Identical networks serve the DC bus ($f_c\approx\SI{1.6}{\kilo\hertz}$) and
the three phase nodes ($f_c\approx$\,2--\SI{3}{\kilo\hertz}).}
\label{fig:divider}
\end{figure}

%=====================================================================
\section{RC Snubber on the Switching Node}\label{sec:snubber}

\subsection{Origin of the ringing}
At each turn-off, the parasitic inductance $L_p$ of the commutation loop (package leads, PCB traces,
bus-capacitor ESL) resonates with the MOSFET output capacitance $C_o$ ($C_{oss}$ plus mounting
capacitance), producing a damped oscillation on the switching node at
\begin{equation}
f_{\mathrm{ring}}=\frac{1}{2\pi\sqrt{L_p C_o}},
\label{eq:fring}
\end{equation}
typically in the tens of MHz. The ringing stresses the device voltage, radiates EMI, and couples
into signal lines. An RC damping snubber placed across the device adds a resistor $R_s$ (in series
with a capacitor $C_s$) that dissipates the resonant energy~\cite{severns}.

\subsection{Design methodology}
The rigorous, layout-independent procedure measures the ringing on the assembled board~\cite{severns}:
\begin{enumerate}[leftmargin=1.4em]
\item Measure the ring period without a snubber, $T_1=1/f_{\mathrm{ring}}$, on the switching node.
\item Add a known capacitor $C_{\mathrm{test}}$ across the device and measure the new period $T_2$.
\item Since $T_1=2\pi\sqrt{L_pC_o}$ and $T_2=2\pi\sqrt{L_p(C_o+C_{\mathrm{test}})}$, squaring and
subtracting eliminates $L_p$:
\begin{equation}
L_p=\frac{T_2^{2}-T_1^{2}}{4\pi^{2}\,C_{\mathrm{test}}},\qquad
C_o=\frac{T_1^{2}}{4\pi^{2}L_p}.
\label{eq:extract}
\end{equation}
\end{enumerate}
The damping resistor is set to the characteristic impedance of the resonance,
\begin{equation}
R_s=\sqrt{\frac{L_p}{C_o}},
\label{eq:rs}
\end{equation}
giving approximately critical damping~\cite{severns}. The snubber capacitor is chosen a few times
larger than $C_o$ so that $R_s$ dominates the loop, yet small enough to limit dissipation; the
resistor power is the energy stored in $C_s$, charged and discharged once per switching cycle,
\begin{equation}
P_{R_s}=C_s\,\Vbus^{2}\,\fsw .
\label{eq:psnub}
\end{equation}

\subsection{Estimated values (to be validated on the bench)}
No assembled board exists yet, so $L_p$ cannot be measured; the values below are a first estimate for
footprint selection only and \emph{must} be confirmed with the two-measurement method of
\eqref{eq:extract} on the prototype. From the IPT015N10N5 datasheet, $C_{oss}=\SI{1800}{\pico\farad}$
(at the datasheet test condition; $C_{oss}$ is strongly voltage-dependent)~\cite{ipt015ds,ipt015part};
adding $\approx\SI{200}{\pico\farad}$ of mounting capacitance gives $C_o\approx\SI{2}{\nano\farad}$.
Taking a representative $L_p\approx\SI{10}{\nano\henry}$ for a compact layout,
\begin{align}
f_{\mathrm{ring}}&=\frac{1}{2\pi\sqrt{(10\text{n})(2\text{n})}}=\SI{35.6}{\mega\hertz},\\
R_s&=\sqrt{\frac{\SI{10}{\nano\henry}}{\SI{2}{\nano\farad}}}=\SI{2.24}{\ohm}\;\to\;\SI{2.2}{\ohm}\ (\text{E12}).
\end{align}
Choosing $C_s=\SI{2.2}{\nano\farad}$ ($\approx C_o$), the worst-case dissipation (\SI{60}{\volt} bus)
from \eqref{eq:psnub} is
\begin{equation}
P_{R_s}=\SI{2.2}{\nano\farad}\times 60^{2}\times\SI{40}{\kilo\hertz}=\SI{0.32}{\watt}\ \text{per snubber},
\end{equation}
i.e.\ $\le\SI{1.9}{\watt}$ for all six switches if every footprint is populated. A
$\ge\SI{1}{\watt}$ resistor is specified for margin.

\subsection{Component selection}
$C_s$ must be C0G/NP0: Class-2 (X7R) ceramics lose much of their capacitance under DC bias and
self-heat under the high-frequency ripple, whereas C0G is lossless and stable. A
\SI{2.2}{\nano\farad}, \SI{100}{\volt} C0G part is the baseline (a \SI{200}{\volt} part adds reserve
against ringing overshoot). $R_s$ must be non-inductive: an SMD thin/thick-film chip (1206 or 2512),
never wirewound, so that its own inductance does not defeat the snubber~\cite{severns}. The RC is
placed directly across the drain--source of each switch with the shortest possible loop. Per the
specification, all six RC sites are laid out as \textbf{do-not-populate (DNP)} footprints: if the
prototype shows acceptable ringing they remain empty, and if not, they are populated using the
bench-measured values from \eqref{eq:extract}--\eqref{eq:rs}. Fig.~\ref{fig:snub} shows the placement
and qualitative effect.

\begin{figure}[!t]
\centering
\begin{circuitikz}[scale=0.95,transform shape,american]
\draw (0,3.4) node[anchor=east]{switch node} -- (1.2,3.4);
\draw (1.2,2.6) rectangle (2.2,3.4);
\draw (1.7,3.0) node{\small $Q$};
\draw (1.7,2.6) -- (1.7,1.6) node[ground]{};
\draw (1.7,3.4) -- (3.0,3.4) to[R,l=$R_s$\,\SI{2.2}{\ohm}] (3.0,2.6)
      to[C,l=$C_s$\,\SI{2.2}{\nano\farad}] (3.0,1.6) node[ground]{};
\draw (1.7,3.4) node[circ]{};
\draw (3.0,3.4) node[anchor=south]{\scriptsize DNP RC};
\begin{scope}[shift={(4.6,1.7)}]
\draw[->] (0,0) -- (3.2,0) node[anchor=north]{\scriptsize $t$};
\draw[->] (0,-0.9) -- (0,1.5) node[anchor=east]{\scriptsize $v$};
\draw[blue,thick] plot coordinates {(0.00,0.85) (0.05,0.73) (0.11,0.46) (0.16,0.10) (0.21,-0.26) (0.27,-0.54) (0.32,-0.68) (0.38,-0.66) (0.43,-0.50) (0.48,-0.24) (0.54,0.06) (0.59,0.33) (0.64,0.50) (0.70,0.56) (0.75,0.49) (0.80,0.31) (0.86,0.08) (0.91,-0.16) (0.96,-0.34) (1.02,-0.44) (1.07,-0.44) (1.12,-0.34) (1.18,-0.17) (1.23,0.03) (1.29,0.20) (1.34,0.33) (1.39,0.37) (1.45,0.33) (1.50,0.22) (1.55,0.06) (1.61,-0.09) (1.66,-0.22) (1.71,-0.29) (1.77,-0.29) (1.82,-0.23) (1.88,-0.12) (1.93,0.01) (1.98,0.13) (2.04,0.21) (2.09,0.24) (2.14,0.22) (2.20,0.15) (2.25,0.05) (2.30,-0.05) (2.36,-0.14) (2.41,-0.19) (2.46,-0.19) (2.52,-0.15) (2.57,-0.08) (2.62,0.00) (2.68,0.08) (2.73,0.14) (2.79,0.16) (2.84,0.15) (2.89,0.10) (2.95,0.04) (3.00,-0.03)};
\draw[red,thick] plot coordinates {(0.00,0.85) (0.05,0.67) (0.11,0.38) (0.16,0.08) (0.21,-0.18) (0.27,-0.35) (0.32,-0.41) (0.38,-0.36) (0.43,-0.25) (0.48,-0.11) (0.54,0.03) (0.59,0.13) (0.64,0.18) (0.70,0.18) (0.75,0.15) (0.80,0.09) (0.86,0.02) (0.91,-0.04) (0.96,-0.07) (1.02,-0.09) (1.07,-0.08) (1.12,-0.06) (1.18,-0.03) (1.23,0.00) (1.29,0.03) (1.34,0.04) (1.39,0.04) (1.45,0.03) (1.50,0.02) (1.55,0.01) (1.61,-0.01) (1.66,-0.02) (1.71,-0.02) (1.77,-0.02) (1.82,-0.01) (1.88,-0.01) (1.93,0.00) (2.04,0.01) (2.20,0.00) (2.62,0.00) (3.00,0.00)};
\node[blue,anchor=west] at (1.3,1.15) {\scriptsize without snubber};
\node[red,anchor=west] at (1.3,0.62) {\scriptsize with snubber};
\end{scope}
\end{circuitikz}
\caption{DNP RC snubber across each switch and its qualitative effect on switching-node ringing: the
damping resistor $R_s=\sqrt{L_p/C_o}$ removes the resonant energy. The values shown are first
estimates to be replaced by bench measurements via \eqref{eq:extract}.}
\label{fig:snub}
\end{figure}

%=====================================================================
\section{Grounding, Noise, and Filtering Strategy}\label{sec:noise}

The system uses a \emph{single shared ground}: the high-current power return, the analog/signal
reference, and all external interfaces sit on one common ground net. There is no separate analog
ground plane; at \SI{40}{\kilo\hertz} and currents above \SI{100}{\ampere}, noise immunity is
obtained through layout discipline rather than ground partitioning. Filtering is applied at each
point where a sensitive analog signal is generated or where switching noise is injected:
\begin{itemize}[leftmargin=1.3em]
\item \textbf{CSA output} (Section~\ref{sec:csafilter}): a light \SI{723}{\kilo\hertz} RC
      (\SI{100}{\ohm}/\SI{2.2}{\nano\farad}) per phase, relying on synchronous sampling rather than a
      delay-inducing anti-alias filter~\cite{an5346}.
\item \textbf{CSA input}: the input series resistors must be \emph{matched} between IN+ and IN-- to
      preserve CMRR; the INA240 enhanced PWM rejection is the primary defense against the common-mode
      transient, so a heavy input filter is deliberately avoided~\cite{ina240ds}.
\item \textbf{Bus divider}: \SI{1.6}{\kilo\hertz} RC (Section~\ref{sec:vfilter}).
\item \textbf{Phase dividers}: 2--\SI{3}{\kilo\hertz} RC with firmware phase-delay compensation where
      used for back-EMF observation, plus Schottky clamps.
\item \textbf{Analog supply}: a ferrite bead between $V_{\mathrm{DD}}$ and $V_{\mathrm{DDA}}$ with
      \SI{1}{\micro\farad}\,$+$\,\SI{100}{\nano\farad} decoupling on $V_{\mathrm{DDA}}$; separate
      \SI{1}{\micro\farad}\,$+$\,\SI{100}{\nano\farad} on $V_{\mathrm{REF+}}$, plus the VREFBUF
      decoupling specified by ST if the internal reference is used~\cite{stm32g473ds}.
\item \textbf{INA240 supply}: \SI{100}{\nano\farad}\,$+$\,\SI{1}{\micro\farad} at the pin.
\item \textbf{DC-link decoupling}: low-ESL ceramics in parallel with the bulk electrolytics, close to
      each half-bridge, to minimize the commutation-loop inductance $L_p$ --- attacking the
      \emph{source} of the ringing the snubber would otherwise have to damp~\cite{severns}.
\item \textbf{Layout}: Kelvin (four-wire) sensing across each shunt; short, tightly coupled
      differential sense traces routed away from the high-$\mathrm{d}i/\mathrm{d}t$ switching-current
      return paths; a compact commutation loop; and a solid, low-impedance ground pour so that the
      shared-path voltage drops are minimized. These layout measures are as important as the
      component filters at this power level.
\end{itemize}

%=====================================================================
\section{Thermal Considerations}\label{sec:thermal}

The MOSFETs are mounted to a liquid-cooled cold plate. At the worst-case corner each device carries
$\approx\SI{60}{\ampere}$ RMS, giving a conduction loss of
$\approx 60^2\times\SI{1.5}{\milli\ohm}\approx\SI{5.4}{\watt}$ cold, rising toward $\sim\SI{8}{\watt}$
as $\Rdson$ increases with junction temperature; switching loss at \SI{40}{\kilo\hertz} and the
shunt loss ($\approx\SI{1.4}{\watt}$ per phase at \SI{119}{\ampere} RMS) add to the budget. Using the
device junction-to-case path (\SI{0.4}{\kelvin\per\watt}) and an adequate thermal-interface material
to the cold plate keeps the junction temperature within rating. The coolant inlet temperature and
flow rate govern the junction temperature far more than the air ambient and are set per deployment.
Low-inductance layout and adequate copper for the $\sim\SI{169}{\ampere}$ peak phase current and
$\sim\SI{132}{\ampere}$ bus current remain board-level tasks, essential to realizing the calculated
switching performance.

%=====================================================================
\section{Bill of Materials (Analog Front-End)}\label{sec:bom}

Table~\ref{tab:bom} lists the analog front-end components. The active devices (INA240A2PWR,
IPT015N10N5) and standard chip passives are high-volume parts; the only component to confirm for
availability at order time is the high-power shunt, for which an alternative of two very common
\SI{0.2}{\milli\ohm} 2512 elements in parallel ($=\SI{0.1}{\milli\ohm}$) is given. Quantities assume
one three-phase board; the snubber components are DNP (footprints only).

\begin{table}[!t]
\centering
\caption{Analog front-end bill of materials. Stock varies; confirm at order time.}
\label{tab:bom}
\small
\renewcommand{\arraystretch}{1.2}
\begin{tabularx}{\textwidth}{@{}>{\raggedright\arraybackslash}p{2.7cm} c L >{\raggedright\arraybackslash}p{2.6cm} L@{}}
\toprule
\textbf{Function} & \textbf{Qty} & \textbf{Component / value} & \textbf{Key rating} & \textbf{Note}\\
\midrule
Current-sense amplifier & 3 & TI INA240A2PWR & $G{=}50$, $-4$ to \SI{80}{\volt} CM & one per phase~\cite{lcsc}\\
Phase shunt & 3 & \SI{0.1}{\milli\ohm} metal element, \SI{5}{\watt} & $\le\SI{2}{\nano\henry}$ ESL, Kelvin & confirm stock; see alt.\ below\\
\;(shunt alternative) & 6 & \SI{0.2}{\milli\ohm} 2512 $\times2\,\parallel$ & \SI{2}{\watt} each & very common; $=\SI{0.1}{\milli\ohm}$\\
CSA output filter R & 3 & \SI{100}{\ohm} 0603 1\% & --- & with $C_f$\\
CSA output filter C & 3 & \SI{2.2}{\nano\farad} C0G 0603 & \SI{50}{\volt} & $f_c\,\SI{723}{\kilo\hertz}$\\
Divider top & 4(8) & \SI{110}{\kilo\ohm} ($2\times\SI{55}{\kilo\ohm}$) 0805 1\% & $\ge\SI{150}{\volt}$ & bus $+$ 3 phases\\
Divider bottom & 4 & \SI{4.7}{\kilo\ohm} 0603 1\% & --- & ratio \num{24.4}\\
Bus filter cap & 1 & \SI{22}{\nano\farad} C0G/X7R 0603 & \SI{50}{\volt} & $f_c\,\SI{1.6}{\kilo\hertz}$\\
Phase filter cap & 3 & \SI{22}{\nano\farad} 0603 & \SI{50}{\volt} & $f_c\approx$2--\SI{3}{\kilo\hertz}\\
ADC clamp diode & 4 & BAT54S dual Schottky & \SI{30}{\volt} & to $V_{\mathrm{DDA}}$/GND\\
Snubber R (DNP) & 6 & \SI{2.2}{\ohm} 1206 non-ind., $\ge\SI{1}{\watt}$ & --- & footprint only\\
Snubber C (DNP) & 6 & \SI{2.2}{\nano\farad} C0G 1206 & \SI{100}{\volt} (\SI{200}{\volt} opt.) & footprint only\\
Ferrite bead on VDDA & 1 & \SI{600}{\ohm}@\SI{100}{\mega\hertz} 0805 & --- & $V_{\mathrm{DD}}\!\to\!V_{\mathrm{DDA}}$\\
Decoupling set & --- & \SI{100}{\nano\farad}/\SI{1}{\micro\farad} 0402/0603 & \SI{16}{\volt}+ & CSA, VDDA, VREF\\
Power MOSFET (ref.) & 12 & Infineon IPT015N10N5 & \SI{100}{\volt}/\SI{1.5}{\milli\ohm} & power stage~\cite{lcsc}\\
\bottomrule
\end{tabularx}
\end{table}

%=====================================================================
\section{Conclusion}\label{sec:concl}

A complete power, gate-drive, and analog input stage has been designed and costed for a
\SIrange{19}{60}{\volt}, \mbox{2.5--5\,kW}, \SI{40}{\kilo\hertz} BLDC/PMSM FOC inverter. The Infineon
IPT015N10N5 and Texas Instruments UCC27211A-Q1 form a compatible pair across the full battery range:
at the highest bus voltage (\SI{60}{\volt}) the \SI{100}{\volt} MOSFET gives a $\approx1.67\times$ DC
blocking margin with about \SI{40}{\volt} of reserve for turn-off overshoot, and at the lowest bus
voltage (\SI{19}{\volt}) the worst-case current reaches $\approx\SI{132}{\ampere}$ on the bus
($\approx\SI{60}{\ampere}$ per device), far below the $2\times\SI{250}{\ampere}$ rating of the
paralleled pair. One channel of the \SI{120}{\volt}, \SI{3.7}{\ampere}/\SI{4.5}{\ampere} driver
switches the \emph{paralleled} gate charge in under \SI{30}{\nano\second} at a moderate
$\approx\SI{2}{\volt\per\nano\second}$, far below its \SI{50}{\volt\per\nano\second} limit. The main
consequence of paralleling two high-$Q_g$ devices is that the bootstrap and $V_{DD}$ decoupling
capacitors must be $\approx$1--2\,\si{\micro\farad} rather than the datasheet's nominal
\SI{0.1}{\micro\farad}, a result that follows directly from the charge-droop equations. The complete
gate-drive passive set is $C_{BOOT}=\SI{2.2}{\micro\farad}$, $C_{VDD}=\SI{4.7}{\micro\farad}+\SI{0.1}{\micro\farad}$,
an asymmetric gate network of $R_{on}=\SI{3.3}{\ohm}$ with a parallel $R_{off}=\SI{1.0}{\ohm}$ and OFF
Schottky (Nexperia PMEG4030ER) per device for soft turn-on and fast turn-off, and $R_{GS}=\SI{10}{\kilo\ohm}$
per device.

The worst-case continuous phase current was derived from the supply/power range as
$\approx\SI{119}{\ampere}$ RMS ($\approx\SI{169}{\ampere}$ peak) at the \SI{2.5}{\kilo\watt}/\SI{19}{\volt}
corner, and used to size a Kelvin \SI{0.1}{\milli\ohm}/\SI{5}{\watt} metal-element shunt read by an
INA240A2 ($G{=}50$) with mid-scale reference, giving a $\pm\SI{290}{\ampere}$-peak linear range at
\SI{0.16}{\ampere/LSB} resolution and $<\SI{0.25}{\ampere}$ offset. The in-line topology is justified
against the low-side, high-side, single-shunt, and Hall alternatives by its continuous,
duty-independent observability, consistent with TI TIDA-00913. For the voltage dividers (ratio
\num{24.4}) it was shown \emph{quantitatively} that no external buffer is required: the
\SI{4.5}{\kilo\ohm} source settles to 12-bit accuracy in $<\SI{300}{\nano\second}$, well within the
available sampling window, and the six STM32G473 internal op-amps remain available as a no-cost
fallback. The RC snubber was sized from the published $C_{oss}$ and the standard two-measurement
ringing-extraction method and laid out as a DNP footprint on all six switches. Filtering was
specified at the CSA output, the dividers, the analog supply/reference, and the DC bus, together with
the single-shared-ground and Kelvin-sensing layout measures that are decisive at this power level.
The next step is prototype bring-up: measuring the actual switching-node ringing to finalize the
snubber, verifying the current-loop noise under synchronous sampling, and confirming shunt
availability at BOM release.

%=====================================================================
\begin{thebibliography}{99}
\setlength{\itemsep}{2pt}

\bibitem{ipt015ds} Infineon Technologies, \emph{IPT015N10N5 --- OptiMOS\texttrademark\,5 Power Transistor, 100\,V}, Final Data Sheet, Rev.\,2.4, 2023. [Online]. Available: \url{https://www.infineon.com/assets/row/public/documents/24/49/infineon-ipt015n10n5-datasheet-en.pdf}

\bibitem{ipt015part} Infineon Technologies, \emph{IPT015N10N5 product page (parametric summary)}. [Online]. Available: \url{https://www.infineon.com/part/IPT015N10N5}

\bibitem{ucc27211} Texas Instruments, \emph{UCC27211A-Q1 Automotive 120-V, 3.7-A/4.5-A Half-Bridge Driver with 8-V UVLO}, Datasheet SLUSCG0C, Dec.\,2015 (rev.\,Sep.\,2025). [Online]. Available: \url{https://www.ti.com/lit/ds/symlink/ucc27211a-q1.pdf}

\bibitem{ucc27211folder} Texas Instruments, \emph{UCC27211A-Q1 product folder}. [Online]. Available: \url{https://www.ti.com/product/UCC27211A-Q1}

\bibitem{balogh} L.\,Balogh, \emph{Fundamentals of MOSFET and IGBT Gate Driver Circuits}, Texas Instruments Application Report SLUA618A, 2017 (rev.\,2018). [Online]. Available: \url{https://www.ti.com/lit/pdf/slua618}

\bibitem{infineon_monolithic} Infineon Technologies, \emph{Using Monolithic High-Voltage Gate Drivers}, Application Note. [Online]. Available: \url{https://www.infineon.com/assets/row/public/documents/24/42/infineon-using-monolithic-high-voltage-gate-drivers-applicationnotes-en.pdf}

\bibitem{ir_an978} International Rectifier (now Infineon Technologies), \emph{Application Note AN-978: HV Floating MOS-Gate Driver ICs}. [Online]. Available: \url{https://www.infineon.com/assets/row/public/documents/24/42/infineon-hv-floating-mos-gate-drivers-applicationnotes-en.pdf}

\bibitem{ir_an1123} A.\,Merello, International Rectifier (now Infineon Technologies), \emph{Application Note AN-1123: Bootstrap Network Analysis --- Focusing on the Integrated Bootstrap Functionality}. [Online]. Available: \url{https://www.infineon.com/assets/row/public/documents/cross-divisions/42/infineon-bootstrap-network-analysis-applicationnotes-en.pdf}

\bibitem{pmeg4030er} Nexperia, \emph{PMEG4030ER --- 40\,V, 3\,A low-$V_F$ Schottky barrier rectifier}, Data Sheet, 2023. [Online]. Available: \url{https://assets.nexperia.com/documents/data-sheet/PMEG4030ER.pdf}

\bibitem{an90059} Nexperia, \emph{Power MOSFET Gate Driver Fundamentals}, Application Note AN90059, Rev.\,1, 2025. [Online]. Available: \url{https://assets.nexperia.com/documents/application-note/AN90059.pdf}

\bibitem{toshiba_dvdt} Toshiba Electronic Devices \& Storage, \emph{Impacts of the dv/dt Rate on MOSFETs}, Application Note, 2018. [Online]. Available: \url{https://toshiba.semicon-storage.com/info/docget.jsp?did=59464}

\bibitem{ti_foc} Texas Instruments, \emph{Sensored Field Oriented Control of 3-Phase Permanent Magnet Synchronous Motors}, Application Report SPRABZ0A, 2016 (rev.\,2021). [Online]. Available: \url{https://www.ti.com/lit/pdf/sprabz0}

\bibitem{ti_svpwm} Z.\,Yu, \emph{Space-Vector PWM With TMS320C24x/F24x Using Hardware and Software Determined Switching Patterns}, Texas Instruments Application Report SPRA524, 1999. [Online]. Available: \url{https://www.ti.com/lit/pdf/spra524}

\bibitem{stm32g473ds} STMicroelectronics, \emph{STM32G473xB/xC/xE Datasheet}, DS12712, Rev.\,5, 2023. [Online]. Available: \url{https://www.st.com/resource/en/datasheet/stm32g473rc.pdf}

\bibitem{ina240ds} Texas Instruments, \emph{INA240 --4-V to 80-V, Bidirectional, Ultra-Precise Current-Sense Amplifier with Enhanced PWM Rejection}, Datasheet SBOS662C, 2016 (rev.\,2021). [Online]. Available: \url{https://www.ti.com/lit/gpn/INA240}

\bibitem{tida00913} Texas Instruments, \emph{TIDA-00913: 48-V/10-A Three-Phase Inverter with In-Line Shunt-Based Phase Current Sensing Reference Design}. [Online]. Available: \url{https://www.ti.com/tool/TIDA-00913}

\bibitem{ti_inline} Texas Instruments, \emph{Low-Drift, Precision, In-Line Motor Current Measurements with PWM Rejection} (INA240 application material). [Online]. Available: \url{https://www.ti.com/product/INA240}

\bibitem{an5346} STMicroelectronics, \emph{AN5346: STM32G4 ADC Use Tips and Recommendations}, Rev.\,2. [Online]. Available: \url{https://www.st.com/resource/en/application_note/an5346-stm32g4-adc-use-tips-and-recommendations-stmicroelectronics.pdf}

\bibitem{an5306} STMicroelectronics, \emph{AN5306: Operational Amplifier (OPAMP) Usage in STM32G4 Series}. [Online]. Available: \url{https://www.st.com/resource/en/application_note/an5306-operational-amplifier-opamp-usage-in-stm32g4-series-stmicroelectronics.pdf}

\bibitem{an2834} STMicroelectronics, \emph{AN2834: How to Get the Best ADC Accuracy in STM32 Microcontrollers}, Rev.\,10. [Online]. Available: \url{https://www.st.com/resource/en/application_note/an2834-how-to-get-the-best-adc-accuracy-in-stm32-microcontrollers-stmicroelectronics.pdf}

\bibitem{severns} R.\,Severns, \emph{Design of Snubbers for Power Circuits}, Cornell Dubilier Application Note. [Online]. Available: \url{https://www.cde.com/resources/technical-papers/design.pdf}

\bibitem{lcsc} LCSC Electronics, component sourcing: INA240A2PWR (C129949), IPT015N10N5 (C108964). [Online]. Available: \url{https://www.lcsc.com/}

\end{thebibliography}

\end{document}
